% Problems, Chapter 22: Calculus of Variations
% Arfken, Weber & Harris, Mathematical Methods for Physicists, 7th ed.
% Requires: amsmath, amssymb, tikz

\section*{Problems}

\Needspace{8\baselineskip}
\subsection*{22.1 Euler Equation}

\Needspace{5\baselineskip}
\par\bigskip
\noindent\textbf{22.1.1}\quad For $dy/dx\equiv y_{x}\ne0$, show
the equivalence of the two forms of Euler's equation:
\[
\erratum{\frac{\partial f}{\partial y}}-\frac{d}{dx}\,
\frac{\partial f}{\partial y_{x}}=0
\]
and
\[
\erratum{\frac{\partial f}{\partial x}}-\frac{d}{dx}
\left(f-y_{x}\frac{\partial f}{\partial y_{x}}\right)=0.
\]
\solution
\erratum{The two partial derivatives in the original textbook statements are
interchanged. The correct equivalent forms are}
\[
\boxed{\frac{\partial f}{\partial y}
-\frac{d}{dx}\frac{\partial f}{\partial y_x}=0},\qquad
\boxed{\frac{\partial f}{\partial x}
-\frac{d}{dx}\left(f-y_x\frac{\partial f}{\partial y_x}\right)=0}.
\]
To prove the equivalence, put $p=y_x$, so $p'=y_{xx}$, and regard
$f=f(y,p,x)$. Subscripts denote its partial derivatives, evaluated on
the curve $y(x)$; $d/dx$ is the total derivative along the curve.
The chain rule gives
\[
\frac{df}{dx}=f_x+f_yp+f_p p'.
\]
By the product rule, $d(pf_p)/dx=p'f_p+p\,df_p/dx$. Therefore
\[
\begin{aligned}
\frac{d}{dx}(f-pf_p)
&=f_x+f_yp+f_p p'-p'f_p-p\frac{df_p}{dx}\\
&=f_x+p\left(f_y-\frac{df_p}{dx}\right),
\end{aligned}
\]
and consequently
\[
f_x-\frac{d}{dx}(f-pf_p)
=-p\left(f_y-\frac{df_p}{dx}\right).
\]
The second corrected expression is $-p$ times the first. For $p\ne0$,
either one vanishes exactly when the other does. At $p=0$ the second
can vanish without implying the first: for $f=y^2$ and $y\equiv1$,
$f_x-d(f-pf_p)/dx=-d(1)/dx=0$, whereas $f_y-df_p/dx=2$.

The printed versions are not generally equivalent. For $f=y$ and
$y(x)=2x$ we have $f_x=f_p=0$, $f_y=1$ and $f-pf_p=y=2x$, so the first
printed expression is $f_x-df_p/dx=0$,
whereas the second is $f_y-d(f-pf_p)/dx=1-2=-1$.
This also shows why the stated nonzero-slope condition does not repair
the swapped derivatives.


\Needspace{5\baselineskip}
\par\bigskip
\noindent\textbf{22.1.2}\quad Derive Euler's equation by expanding
the integrand of
\[
J(\alpha)=\int_{x_{1}}^{x_{2}}
f\big(y(x,\alpha),y_{x}(x,\alpha),x\big)\,dx
\]
in powers of $\alpha$.
\erratum{Keep both endpoint values fixed and assume $f$ is $C^3$
near the smooth trial curves, so that the expansion through order
$\alpha^2$ has an $O(\alpha^3)$ remainder.}
\par\emph{Note.}~The stationary condition is
$\partial J(\alpha)/\partial\alpha=0$, evaluated at $\alpha=0$.
The terms quadratic in $\alpha$ may be useful in establishing the
nature of the stationary solution (maximum, minimum, or saddle
point).
\solution
Choose $y(x,\alpha)=y(x)+\alpha\eta(x)$ with
$\eta(x_1)=\eta(x_2)=0$, so
$y_x(x,\alpha)=y'(x)+\alpha\eta'(x)$. Taylor expansion about
$\alpha=0$ gives
\[
\begin{aligned}
f(y+\alpha\eta,y'+\alpha\eta',x)
={}&f+\alpha(f_y\eta+f_{y_x}\eta')\\
&+\frac{\alpha^2}{2}
\left(f_{yy}\eta^2+2f_{y y_x}\eta\eta'
+f_{y_xy_x}(\eta')^2\right)+O(\alpha^3).
\end{aligned}
\]
All derivatives on the right are evaluated on the unperturbed curve.
Integrating over $[x_1,x_2]$ (the remainder is $O(\alpha^3)$ uniformly
in $x$ because $f$ is $C^3$ and $\eta,\eta'$ are bounded) gives
\[
J(\alpha)=J(0)+\alpha J'(0)+\frac{\alpha^2}{2}J''(0)+O(\alpha^3),
\]
where $J'(0)$ and $J''(0)$ are the integrals of the coefficients of
$\alpha$ and $\alpha^2/2$ above. Integrating the linear term by parts,
$\int f_{y_x}\eta'dx=[f_{y_x}\eta]-\int(df_{y_x}/dx)\eta\,dx$, so
\[
\begin{aligned}
J'(0)&=\int_{x_1}^{x_2}(f_y\eta+f_{y_x}\eta')dx\\
&=[f_{y_x}\eta]_{x_1}^{x_2}
+\int_{x_1}^{x_2}\left(f_y-\frac d{dx}f_{y_x}\right)\eta\,dx.
\end{aligned}
\]
The boundary term is zero because $\eta(x_1)=\eta(x_2)=0$.
Stationarity, $J'(0)=0$, must hold for every such $\eta$. The bracket
is continuous; if it were nonzero at an interior point, it would keep
one sign on a small interval there, and a nonnegative $\eta\ne0$
vanishing outside that interval would give $J'(0)\ne0$. This
fundamental lemma of the calculus of variations gives
\[
\boxed{f_y-\frac d{dx}f_{y_x}=0}.
\]
The coefficient of $\alpha^2$ is one half of
\[
J''(0)=\int_{x_1}^{x_2}
\left(f_{yy}\eta^2+2f_{y y_x}\eta\eta'
+f_{y_xy_x}(\eta')^2\right)dx.
\]
If $J''(0)>0$ for every admissible $\eta\ne0$, then $J(\alpha)>J(0)$
for all small $\alpha\ne0$ in every such family (a minimum along each
family); $J''(0)<0$ for every $\eta\ne0$ gives a maximum in the same
sense, and values of both signs give a saddle point. Thus the sign must
be tested over all admissible variations; the vanishing first variation
alone does not classify the extremum.


\Needspace{5\baselineskip}
\par\bigskip
\noindent\textbf{22.1.3}\quad Find the Euler equation
corresponding to Eq.~(22.14) if $f=f(y_{xx},y_{x},y,x)$, assuming
that $y$ and $y_{x}$ have fixed values at the endpoints of their
interval of definition.
\par\nobreak\hfill\emph{ANS.}\enspace
$\dfrac{d^{2}}{dx^{2}}\left(\dfrac{\partial f}{\partial y_{xx}}
\right)-\dfrac{d}{dx}\left(\dfrac{\partial f}{\partial y_{x}}
\right)+\dfrac{\partial f}{\partial y}=0$.
\solution
Write $y_\alpha=y+\alpha\eta$, so $y_\alpha'=y'+\alpha\eta'$ and
$y_\alpha''=y''+\alpha\eta''$. Fixing both $y$ and $y'$ at each
endpoint requires $\eta=\eta'=0$ there. Differentiating
$J(\alpha)=\int f(y_\alpha'',y_\alpha',y_\alpha,x)\,dx$ under the
integral sign by the chain rule, the first variation
$\delta J=J'(0)$ is
\[
\delta J=\int_{x_1}^{x_2}
(f_y\eta+f_{y_x}\eta'+f_{y_{xx}}\eta'')dx.
\]
The first derivative term gives
\[
\int f_{y_x}\eta'dx
=[f_{y_x}\eta]-\int\frac d{dx}f_{y_x}\,\eta\,dx.
\]
For the second derivative term, integrate twice:
\[
\int f_{y_{xx}}\eta''dx
=[f_{y_{xx}}\eta']-\int\frac d{dx}f_{y_{xx}}\eta'dx
=[f_{y_{xx}}\eta'-(d f_{y_{xx}}/dx)\eta]
+\int\frac{d^2}{dx^2}f_{y_{xx}}\eta\,dx.
\]
Every displayed boundary term vanishes, since each contains $\eta$ or
$\eta'$ at $x_1$ or $x_2$. Thus
\[
\delta J=\int_{x_1}^{x_2}
\left[f_y-\frac d{dx}f_{y_x}
+\frac{d^2}{dx^2}f_{y_{xx}}\right]\eta\,dx.
\]
This vanishes for every such $\eta$, so the fundamental lemma
(Exercise~22.1.2) yields
\[
\boxed{\frac{d^2}{dx^2}\frac{\partial f}{\partial y_{xx}}
-\frac d{dx}\frac{\partial f}{\partial y_x}
+\frac{\partial f}{\partial y}=0}.
\]


\Needspace{5\baselineskip}
\par\bigskip
\noindent\textbf{22.1.4}\quad The integrand $f(y,y_{x},x)$ of
Eq.~(22.2) has the form
\[
f(y,y_{x},x)=f_{1}(x,y)+f_{2}(x,y)y_{x}.
\]
\begin{enumerate}
\item[(a)] Show that the Euler equation leads to
\[
\frac{\partial f_{1}}{\partial y}-\frac{\partial f_{2}}{\partial x}
=0.
\]
\item[(b)] \erratum{For $C^1$ functions $f_1,f_2$, distinguish this
equality along one stationary curve from an identity on a simply
connected region. Under which condition is $J$ independent of the path?}
\end{enumerate}
\solution
(a) For $f=f_1(x,y)+f_2(x,y)y'$,
\[
f_y=f_{1y}+f_{2y}y',\qquad
f_{y'}=f_2,\qquad
\frac d{dx}f_{y'}=f_{2x}+f_{2y}y'.
\]
The $y'$ terms cancel in Euler's equation, leaving
\[
\boxed{f_{1y}-f_{2x}=0\quad\hbox{along a stationary curve}}.
\]

(b) Since $dy=y'dx$ along the path, the integral is the line integral
\[
J=\int(f_1\,dx+f_2\,dy).
\]
If the equality $f_{1y}=f_{2x}$ holds identically on a simply
connected region, this differential form is exact. Indeed, by Green's
theorem
$\oint(f_1\,dx+f_2\,dy)=\iint(f_{2x}-f_{1y})\,dx\,dy=0$ around every
simple closed curve in the region (its interior lies in the region),
so $F(x,y)=\int_{(x_1,y_1)}^{(x,y)}(f_1\,dx+f_2\,dy)$ does not depend
on the path of integration, and $F_x=f_1$, $F_y=f_2$. Along any path
$y(x)$, $dF(x,y(x))/dx=F_x+F_yy'=f_1+f_2y'=f$. Then
\[
\boxed{J=F(x_2,y_2)-F(x_1,y_1)}
\]
is independent of the path in that region. Conversely, if $J$ has the
same value on all paths joining any two points of the region, every
path is stationary, and (a) gives $f_{1y}=f_{2x}$ along every curve,
hence identically. So $J$ is path-independent exactly when
$f_{1y}=f_{2x}$ holds throughout the (simply connected) region.

\erratum{An Euler equation satisfied only on one curve does not
imply path independence throughout the region. The identity on the
region, together with an appropriate topology, is needed for that
conclusion.}
For example, $f_1=y^2$, $f_2=0$ gives $f_{1y}-f_{2x}=2y$, which
vanishes only on the stationary curve $y=0$. With both endpoints on the
axis, $\int y^2dx$ is $0$ on $y=0$ but positive on every other path,
so it changes when the interior path changes. This example
also anticipates the correction required in the next exercise.


\Needspace{5\baselineskip}
\par\bigskip
\noindent\textbf{22.1.5}\quad Show that the condition that
$J=\displaystyle\int f(x,y)\,dx$ has a stationary value
\begin{enumerate}
\item[(a)] \erratum{leads to $f_y(x,y(x))=0$ along the stationary curve and}
\item[(b)] \erratum{imposes no separate condition $f_x=0$.}
\end{enumerate}
\erratum{Dependence on derivatives of $y$ is not essential for a
meaningful variational problem. Give a smooth minimizing solution
for an integrand that depends on $y$.}
\par\emph{Note.}~The situation will change when constraints are
introduced (compare to Exercise~22.4.6).
\solution
There is no $y'$ dependence, so for $y+\alpha\eta$ with
$\eta(x_1)=\eta(x_2)=0$ the first variation is simply
\[
\delta J=\int_{x_1}^{x_2}f_y(x,y(x))\eta(x)\,dx.
\]
It must vanish for every such $\eta$, so by the fundamental lemma
(Exercise~22.1.2), or by Euler's equation with $f_{y'}=0$, the correct
stationarity condition is the algebraic equation
\[
\boxed{f_y(x,y(x))=0\quad\hbox{for every interior }x}.
\]
\erratum{This is a condition along the selected curve. It does not
force $f(x,y)$ to be independent of $y$, and absence of derivative
dependence does not rule out meaningful, smooth minimizing solutions.
The printed part (a), which concludes that $f(x,y)$ is independent of
$y$, and the printed remark that there is then no continuous,
differentiable solution, are both false.}

For an explicit smooth counterexample take
\[
f(x,y)=(y-\sin x)^2
\]
and fix the endpoint values to those of $\sin x$. Then
$f_y=2(y-\sin x)$ vanishes only for $y=\sin x$, which also meets the
endpoint values, so the stationary function is
$\boxed{y(x)=\sin x}$. Moreover,
\[
J[y]=\int_{x_1}^{x_2}(y-\sin x)^2dx\ge0,\qquad
J[\sin x]=0,
\]
which proves a global minimum despite explicit dependence on $y$.

If $f$ is independent of $y$, every path has the same $J$.
If $f_y$ never vanishes, no stationary path exists. Otherwise, smooth
root branches of $f_y=0$ may furnish stationary paths, subject to the
endpoint conditions. Regarding (b), there is no additional equation
$f_x=0$: $x$ is the independent variable and only $y$ is varied. The
second form of Euler's equation (Exercise~22.1.1) adds nothing either:
with $f_{y'}=0$ it reads
\[
f_x-\frac{df}{dx}=f_x-(f_x+f_yy')=-f_yy'=0,
\]
which already follows from $f_y=0$. But $x$ can affect the
selected curve through $f_y(x,y)=0$, as the counterexample demonstrates.


\Needspace{5\baselineskip}
\par\bigskip
\noindent\textbf{22.1.6}\quad A soap film stretched between two
rings of unit radius centered at $\pm x_{0}$ will have its closest
approach to the $x$-axis at $x=0$, with the distance from the
axis given by $c_{1}$, with $x_{0}$ and $c_{1}$ related by
Eq.~(22.26) or Eq.~(22.29).
\begin{enumerate}
\item[(a)] Show that $dc_{1}/dx_{0}$ becomes infinite when
$x_{0}\sinh(x_{0}/c_{1})=1$, indicating that the soap film
becomes unstable if $x_{0}$ is increased beyond the value
satisfying this condition.
\item[(b)] Show that the condition of part (a) is equivalent to
\[
\frac{x_{0}}{c_{1}}=\coth\frac{x_{0}}{c_{1}}.
\]
\item[(c)] Solve the transcendental equation of part (b) to
obtain the critical value of $x_{0}/c_{1}$ and show that the
separate values of $x_{0}$ and $c_{1}$ are then approximately
$x_{0}\approx0.6627$ and $c_{1}\approx0.5524$.
\end{enumerate}
\solution
Put $c=c_1>0$ and $u=x_0/c$. The catenoid $y=c\cosh(x/c)$ must have
radius $1$ at $x=\pm x_0$, so the boundary equation is
$c\cosh u=1$.

(a) Write $c'=dc/dx_0$. Since $u=x_0/c$, the quotient rule gives
$du/dx_0=1/c-x_0c'/c^2$. Implicit differentiation of $c\cosh u=1$
with respect to $x_0$ therefore gives
\[
c'\cosh u+c\sinh u
\left(\frac1c-\frac{x_0c'}{c^2}\right)=0,
\]
that is, $c'\cosh u+\sinh u-uc'\sinh u=0$; hence
\[
c'(\cosh u-u\sinh u)=-\sinh u.
\]
Multiply the numerator and denominator of
$c'=-\sinh u/(\cosh u-u\sinh u)$ by $c$ and use $c\cosh u=1$,
$cu=x_0$:
\[
\boxed{\frac{dc}{dx_0}=
-\frac{c\sinh u}{1-x_0\sinh u}}.
\]
The numerator is nonzero for separated rings ($x_0>0$, so $u>0$).
Thus the derivative diverges at $x_0\sinh u=1$.

(b) Divide that condition by $c\cosh u=1$; since $x_0/c=u$,
\[
\frac{x_0\sinh u}{c\cosh u}=u\tanh u=1,\qquad \boxed{u=\coth u}.
\]
Equivalently, the boundary equation parametrizes all solutions by
$u>0$:
\[
c=\operatorname{sech}u,\qquad x_0=cu=u\operatorname{sech}u.
\]
By the product rule,
\[
\frac{dx_0}{du}=\operatorname{sech}u-u\operatorname{sech}u\tanh u
=\operatorname{sech}u\,(1-u\tanh u).
\]
Because $d(u\tanh u)/du=\tanh u+u\operatorname{sech}^2u>0$, the
product $u\tanh u$ increases strictly from zero to infinity. So $x_0(u)$
increases while $u\tanh u<1$ and decreases afterwards: there is one
critical point, a maximum, at the root $u_c$ of $u=\coth u$. There
$dc/dx_0=(dc/du)/(dx_0/du)$ is infinite, since
$dc/du=-\operatorname{sech}u\tanh u\ne0$, in agreement with (a).
Each $x_0$ below the maximum is reached twice: on the shallow branch
($u<u_c$, larger neck radius $c$) and on the deep branch ($u>u_c$).
Beyond the maximum no positive-neck catenoid satisfies the two boundary
conditions.

(c) Newton iteration for $F(u)=u\tanh u-1$, with
$F'(u)=\tanh u+u\operatorname{sech}^2u$, is
\[
u_{j+1}=u_j-
\frac{u_j\tanh u_j-1}{\tanh u_j+u_j\operatorname{sech}^2u_j}.
\]
Starting from the guess $1.2$, two steps already give all the digits
shown below; then $c_c=\operatorname{sech}u_c$ (from $c\cosh u=1$) and
$x_{0c}=u_cc_c$:
\[
\boxed{u_c=1.19967864025773,\quad
c_c=0.552434124530883,\quad
x_{0c}=0.662743419349182}.
\]
\erratum{The Instructor's Manual's $c_c=0.5523434$ is a misprint;
substituting $c_c=\operatorname{sech}u_c$ gives the value above, which
agrees with the textbook's rounded $0.5524$.}

The stability threshold can also be seen in the second variation
$J''(0)$ of Exercise~22.1.2, written $\delta^2J$, for the area integrand
without its factor $2\pi$, $f=y\sqrt{1+y'^2}$. Its second partial
derivatives with respect to $y$ and $y'$ are
\[
 f_{yy}=0,\qquad f_{yy'}=\frac{y'}{\sqrt{1+y'^2}},\qquad
 f_{y'y'}=\frac y{(1+y'^2)^{3/2}}.
\]
On the catenoid, $y=c\cosh(x/c)$, $y'=\sinh(x/c)$ and
$\sqrt{1+y'^2}=\cosh(x/c)$, so $f_{yy'}=\tanh(x/c)$ and
$f_{y'y'}=c\operatorname{sech}^2(x/c)$.
Thus, for a fixed-endpoint variation $\eta$,
\[
 \delta^2J=\int_{-x_0}^{x_0}
 \left[2\tanh(x/c)\eta\eta_x
       +c\operatorname{sech}^2(x/c)\eta_x^2\right]dx.
\]
The first term is $\tanh(x/c)(\eta^2)_x$. Integration by parts
has zero endpoint term, because $\eta(\pm x_0)=0$, and leaves
$-\int[d\tanh(x/c)/dx]\eta^2dx
=-c^{-1}\int\operatorname{sech}^2(x/c)\eta^2dx$.
Now let $u=x/c$ run over $[-u_0,u_0]$, where $u_0=x_0/c$ is the
endpoint value called $u$ above. Then $dx=c\,du$ and
$\eta_x=\eta_u/c$ turn
$c\operatorname{sech}^2(x/c)\eta_x^2dx$ into
$\operatorname{sech}^2u\,\eta_u^2du$ and
$c^{-1}\operatorname{sech}^2(x/c)\eta^2dx$ into
$\operatorname{sech}^2u\,\eta^2du$, so
\[
\delta^2J=\int_{-u_0}^{u_0}
\operatorname{sech}^2u\,(\eta_u^2-\eta^2)du.
\]
Put $p(u)=\operatorname{sech}^2u$ and
\[
w(u)=\cosh u-u\sinh u=\cosh u\,(1-u\tanh u),
\]
which is $\partial y/\partial c=\cosh(x/c)-(x/c)\sinh(x/c)$ for the
family $y=c\cosh(x/c)$ at fixed $x$, written in terms of $u=x/c$.
Since $w'=-u\cosh u$, we have $pw'=-u\operatorname{sech}u$ and
\[
(pw')'=-\operatorname{sech}u+u\operatorname{sech}u\tanh u=-pw,
\]
that is, $(pw')'+pw=0$. For $u_0<u_c$ (the shallow branch),
$|u|\tanh|u|<1$ on $[-u_0,u_0]$, so $w>0$ there. Write $\eta=wv$;
then $v=\eta/w$ is smooth, $v(\pm u_0)=0$, and $\eta'=w'v+wv'$.
Expanding the square,
\[
 p(\eta'^2-\eta^2)
 =pw^2v'^2+2pww'vv'+(pw'^2-pw^2)v^2
 =pw^2v'^2+(pww'v^2)',
\]
because $(pww'v^2)'=(pw')'wv^2+pw'^2v^2+2pww'vv'$ and
$(pw')'w=-pw^2$. The boundary term $[pww'v^2]_{-u_0}^{u_0}$ is zero
because $v(\pm u_0)=0$. Consequently
\[
\delta^2J=\int_{-u_0}^{u_0}
\operatorname{sech}^2u\,w^2(v')^2du>0
\]
for a nonzero fixed-endpoint variation: $v'\equiv0$ together with
$v(\pm u_0)=0$ would give $v\equiv0$, hence $\eta\equiv0$. At
$u_0=u_c$, $w(\pm u_c)=\cosh u_c(1-u_c\tanh u_c)=0$, so $\eta=w$ is
itself an admissible variation, and integration by parts with
$(pw')'=-pw$ gives
$\delta^2J[w]=[pww']_{-u_c}^{u_c}-\int w[(pw')'+pw]\,du=0$.
This identifies the loss of strict stability on the shallow branch.


\Needspace{5\baselineskip}
\par\bigskip
\noindent\textbf{22.1.7}\quad A soap film is stretched across the
space between two rings of unit radius centered at $\pm x_{0}$ on
the $x$-axis and perpendicular to the $x$-axis. Using the
solution developed in Example~22.1.3, set up the transcendental
equations for the condition that $x_{0}$ is such that the area of
the curved surface of rotation equals the area of the two rings
(Goldschmidt discontinuous solution). Solve for $x_{0}$.
\solution
For $y=c\cosh(x/c)$ between $x=\pm x_0$, put $u=x_0/c$.
The endpoint condition is $c\cosh u=1$.
Since $y'=\sinh(x/c)$, $\sqrt{1+y'^2}=\cosh(x/c)$ and
$y\sqrt{1+y'^2}=c\cosh^2(x/c)$. The integrand is even; with $v=x/c$,
$dx=c\,dv$, and $\cosh^2v=(1+\cosh2v)/2$, the full surface area is
\[
\begin{aligned}
A&=2\pi\int_{-x_0}^{x_0}
y\sqrt{1+y'^2}dx
=4\pi c^2\int_0^u\cosh^2v\,dv\\
&=2\pi c^2\left[v+\tfrac12\sinh2v\right]_0^u
=\pi c^2[\sinh(2u)+2u].
\end{aligned}
\]
The two unit disks have area $2\pi$. Equating the areas, dividing by
$\pi$, and using $c^2=\operatorname{sech}^2u$ gives
\[
\sinh(2u)+2u=2\cosh^2u=1+\cosh(2u).
\]
Since $\sinh2u-\cosh2u=-e^{-2u}$, this is the simpler single equation
\[
\boxed{2u-1=e^{-2u}}.
\]
Its left-minus-right side $g(u)=2u-1-e^{-2u}$ has $g(0)=-2$,
$g(u)\to\infty$ as $u\to\infty$, and $g'(u)=2+2e^{-2u}>0$, so it has
exactly one positive root. Since $g(0.6)=-0.101<0<g(0.7)=0.153$,
bisection on $(0.6,0.7)$, or Newton iteration
$u_{j+1}=u_j-g(u_j)/g'(u_j)$ from the starting value $0.6$, gives
\[
u=0.639232271380537,\qquad
c=\operatorname{sech}u=0.825517453652511
\]
(both sides of the boxed equation then equal $0.278464542761$).
Thus the required half-separation is
\[
\boxed{x_0=cu=u\operatorname{sech}u=0.527697396962572}.
\]
The full ring separation is $2x_0=1.05539479392514$.
The root is below $u_c=1.19967\ldots$ from Exercise~22.1.6, so the
equality occurs on the shallow catenoid branch. With
$c=\operatorname{sech}u$ the area
is $A=2\pi(\tanh u+u\operatorname{sech}^2u)$, and
$dA/du=4\pi\operatorname{sech}^2u\,(1-u\tanh u)>0$ on that branch, as
is $dx_0/du$. Hence the shallow catenoid has less area than the two
disks for $x_0<0.527697\ldots$ and more for
$0.527697\ldots<x_0\le x_{0c}$, where the Goldschmidt solution has less.


\Needspace{5\baselineskip}
\par\bigskip
\noindent\textbf{22.1.8}\quad In Example~22.1.1, expand
$J[y(x,\alpha)]-J[y(x,0)]$ in powers of $\alpha$. The term
linear in $\alpha$ leads to the Euler equation and to the
straight-line solution, Eq.~(22.16). Investigate the $\alpha^{2}$
term and show that the stationary value of $J$, the straight-line
distance, is a minimum.
\solution
Let $y_*(x)=y_1+m(x-x_1)$, where
$m=(y_2-y_1)/(x_2-x_1)$ and $x_2>x_1$.
For $y_\alpha=y_*+\alpha\eta$, with $\eta(x_1)=\eta(x_2)=0$,
\[
J[y_\alpha]=\int_{x_1}^{x_2}
\sqrt{1+(m+\alpha\eta')^2}\,dx.
\]
The first two derivatives of $g(p)=\sqrt{1+p^2}$ are
$g'(p)=p/\sqrt{1+p^2}$ and $g''(p)=(1+p^2)^{-3/2}$.
Taylor's theorem at $p=m$ gives, at each $x$,
\[
g(m+\alpha\eta')=g(m)+g'(m)\,\alpha\eta'
+\tfrac12g''(m)\,\alpha^2(\eta')^2+O(\alpha^3),
\]
with a remainder uniform in $x$, because
$g'''(p)=-3p(1+p^2)^{-5/2}$ and $\eta'$ are bounded. Integrating over
$[x_1,x_2]$,
\[
J[y_\alpha]-J[y_*]
=\alpha\frac{m}{\sqrt{1+m^2}}\int_{x_1}^{x_2}\eta'dx
+\frac{\alpha^2}{2(1+m^2)^{3/2}}
\int_{x_1}^{x_2}(\eta')^2dx+O(\alpha^3).
\]
The linear integral is $\eta(x_2)-\eta(x_1)=0$.
The quadratic coefficient is strictly positive for any nonzero
admissible $\eta$: zero would require $\eta'=0$, and the endpoints
would then require $\eta=0$. Thus the straight line is a strict
local minimum.

It is also the global minimum. Since $g''>0$, $g$ is strictly convex,
so its graph lies above each tangent line: $g(m+t)\ge g(m)+g'(m)t$,
with equality only for $t=0$. With $t=\eta'(x)$,
\[
g(m+\eta')\ge g(m)+g'(m)\eta'.
\]
After integration the last term again vanishes, so
$J[y_*+\eta]\ge J[y_*]$, with equality only if $\eta'\equiv0$, that is,
only for $\eta=0$.


\Needspace{5\baselineskip}
\par\bigskip
\noindent\textbf{22.1.9}\quad
\erratum{Take $x_1<x_2$, real-valued paths, and fixed endpoint values.
In (b), allow endpoint jumps, and compare with the continuous class.}
\begin{enumerate}
\item[(a)] Show that the integral
\[
J=\int_{x_{1}}^{x_{2}}f(y,y_{x},x)\,dx,
\quad\text{with}\quad f=y(x),
\]
has no extreme values.
\item[(b)] If $f(y,y_{x},x)=y^{2}(x)$, find a discontinuous
solution similar to the Goldschmidt solution for the soap-film
problem.
\end{enumerate}
\solution
(a) For $f=y$, Euler's equation $f_y-d(f_{y_x})/dx=0$ reads $1=0$,
so no curve is stationary. Directly: let $\eta$ be a smooth function of
compact support in $(x_1,x_2)$ with $\int\eta\,dx\ne0$. Then
\[
J[y+\alpha\eta]-J[y]=\alpha\int_{x_1}^{x_2}\eta(x)dx.
\]
Arbitrarily small positive and negative $\alpha$ change $J$ in
opposite directions. Thus no admissible curve is even a local
extremum. Taking $|\alpha|\to\infty$ also shows that $J$ is
unbounded above and below.

(b) For $f=y^2$, nonnegativity gives $J\ge0$. Euler's equation is
$f_y=2y=0$ (compare Exercise~22.1.5), which forces $y=0$ at every
interior point; nonzero endpoint values can then be met only by jumps
at the ends, just as the Goldschmidt solution joins the two disks by the
axis. In a class allowing endpoint jumps, set
\[
y(x)=0\quad(x_1<x<x_2),
\]
while assigning the prescribed values at the two endpoints.
Changing isolated points does not change the integral, so this
discontinuous path has $\boxed{J_{\min}=0}$.

\erratum{If continuity is required and either endpoint value is
nonzero, this is an infimum, not an attained minimum. If both endpoint
values are zero, the ordinary smooth function $y=0$ is already the
minimum.}
To see the infimum explicitly, let the endpoint values be $A,B$, and
let $0<\epsilon<(x_2-x_1)/2$. Let $y_\epsilon$ fall linearly from $A$
to zero on $[x_1,x_1+\epsilon]$, vanish on
$[x_1+\epsilon,x_2-\epsilon]$, and rise linearly from zero to $B$ on
$[x_2-\epsilon,x_2]$ (round off the corners if a smooth path is
wanted). Since $|y_\epsilon|\le|A|$ on the first layer and
$|y_\epsilon|\le|B|$ on the last,
\[
0\le J[y_\epsilon]\le\epsilon(A^2+B^2)\longrightarrow0.
\]
On the other hand, a continuous function with $y(x_1)=A\ne0$ satisfies
$|y|>|A|/2$ on some interval $[x_1,x_1+\delta]$, so
$J\ge\delta A^2/4>0$; likewise at $x_2$. Hence no continuous path
attains the value $0$.


\Needspace{5\baselineskip}
\par\bigskip
\noindent\textbf{22.1.10}\quad Fermat's principle of optics states
that a light ray in a medium for which $n$ is the
(position-dependent) index of refraction will follow the path
$y(x)$ for which
\[
\int_{x_{1},y_{1}}^{x_{2},y_{2}}n(y,x)\,ds
\]
is \erratum{stationary}. For $y_{2}=y_{1}=1$, $-x_{1}=x_{2}=1$, find the
ray path \erratum{and determine whether the minimum is attained} if
\begin{enumerate}
\item[(a)] $n=e^{y}$,\qquad
\item[(b)] $n=a(y-y_{0})$,\quad$y>y_{0}$,
\erratum{with $a>0$ and $y_0<1$.}
\end{enumerate}
\solution
Use $ds=\sqrt{1+y'^2}dx$, so $f=n(y)\sqrt{1+y'^2}$ and
$f_{y'}=n(y)y'/\sqrt{1+y'^2}$. Since neither index depends explicitly
on $x$, the Beltrami first integral $f-y'f_{y'}=C$ (the second form of
Euler's equation in Exercise~22.1.1 with $f_x=0$) holds:
\[
f-y'f_{y'}=n(y)\,\frac{(1+y'^2)-y'^2}{\sqrt{1+y'^2}}
=\frac{n(y)}{\sqrt{1+y'^2}}=C.
\]
Conversely, a solution of this first integral satisfies Euler's
equation wherever $y'\ne0$ (Exercise~22.1.1), but not necessarily
where $y'=0$. On an interval where $y$ is constant, $f_{y'}=0$ and
Euler's expression $f_y-df_{y'}/dx$ reduces to $f_y=n'(y)$, which is
$e^y>0$ in (a) and $a>0$ in (b). So a ray contains no horizontal
segment. In particular, the constant path $y\equiv1$ satisfies the
first integral, with $C=n(1)$, and meets both endpoint conditions, but
it is not a ray. The rays found below have $y'=0$ only at $x=0$, so
Euler's equation holds for $x\ne0$ and, by continuity, at $x=0$ as well.

(a) For $n=e^y$ this gives $y'^2=e^{2y}/C^2-1$, with $C>0$. Separate
variables and put $t=Ce^{-y}\in(0,1]$, so that
$\sqrt{e^{2y}/C^2-1}=\sqrt{1-t^2}/t$ and $dt=-t\,dy$:
\[
\pm dx=\frac{dy}{\sqrt{e^{2y}/C^2-1}}
=\frac{t\,dy}{\sqrt{1-t^2}}
=-\frac{dt}{\sqrt{1-t^2}}=d(\arccos t).
\]
Hence $\arccos(Ce^{-y})=\pm(x-b)$, that is, $e^y=C\sec(x-b)$ on a
branch where the cosine is positive. Indeed,
\[
y=\ln C-\ln\cos(x-b),\quad y'=\tan(x-b),
\]
so $\sqrt{1+y'^2}=\sec(x-b)$ and $e^y/\sqrt{1+y'^2}=C$. Equal endpoint
values at $x=\pm1$ require $\cos(1-b)=\cos(1+b)$, i.e.,
$2\sin1\sin b=0$, so $b=k\pi$; since
$\cos(x-k\pi)=(-1)^k\cos x$ must be positive on all of $[-1,1]$, the
branch covering the entire interval has $\cos(x-b)=\cos x$, i.e.,
$b=0$. The condition $y(1)=1$, i.e. $e=C\sec1$, gives $C=e\cos1$. Thus
\[
\boxed{y(x)=1+\ln\cos1-\ln\cos x,\qquad -1\le x\le1}.
\]
On it $e^y=e\cos1\sec x$ and $\sqrt{1+y'^2}=\sec x$, so the optical
length is
\[
J=e\cos1\int_{-1}^1\sec^2x\,dx=2e\cos1\tan1=2e\sin1.
\]
This is a global minimum, not just a stationary path. Set
\[
U=e^y\cos x,\qquad V=e^y\sin x,
\]
polar coordinates with radius $e^y$ and angle $x\in[-1,1]$, so the
map $(x,y)\mapsto(U,V)$ is one-to-one. Since
$dU=e^y(\cos x\,dy-\sin x\,dx)$ and $dV=e^y(\sin x\,dy+\cos x\,dx)$,
the cross terms cancel and $dU^2+dV^2=e^{2y}(dx^2+dy^2)$.
The optical length $\int e^y\,ds$ is therefore ordinary Euclidean
length in the $(U,V)$ plane. The endpoints are $(e\cos1,\pm e\sin1)$,
and their straight segment has $U=e\cos1$, precisely the curve above
(on it $e^y\cos x=C=e\cos1$). Its length is $2e\sin1$, the lower bound
for any joining path; every other path has a different image, which is
strictly longer than the segment.

\erratum{The Instructor's Manual's inverse formula uses $1+\pi^2/16$
where $\sec^2 1$ is required. At $y=1$ it gives
$x=\arctan(\pi/4)$ instead of $x=1$, so it fails the stated
endpoint condition.}

(b) A positive index requires $a>0$ and the endpoints require
$H:=1-y_0>0$. Put $u(x)=y(x)-y_0$, so $n=au$ and $u'=y'$. The first
integral, divided by $a$, is
\[
\frac{u}{\sqrt{1+u'^2}}=c>0,\qquad
u'^2=\frac{u^2}{c^2}-1.
\]
Separate variables on a rising part ($u'>0$), using
$\int du/\sqrt{u^2-c^2}=\operatorname{arcosh}(u/c)$:
\[
dx=\frac{c\,du}{\sqrt{u^2-c^2}},\qquad
x-b=c\,\operatorname{arcosh}(u/c).
\]
A falling part gives $x-b=-c\operatorname{arcosh}(u/c)$; both halves
are $u=c\cosh[(x-b)/c]$. Equal endpoint heights require
$\cosh[(1-b)/c]=\cosh[(1+b)/c]$, i.e. $|1-b|=|1+b|$, so the horizontal
shift is $b=0$:
\[
\boxed{y(x)=y_0+c\cosh(x/c),\qquad
c\cosh(1/c)=H}.
\]
With $w=1/c$, the boundary equation reads $\cosh w/w=H$, or
\[
\frac1H=w\operatorname{sech}w.
\]
Exercise~22.1.6 shows that the maximum of the right side is
$0.662743419349182$, at $w=u_c$. Consequently regular stationary curves
exist only for
\[
H\ge\frac1{0.662743419349182}=1.50887956153832,
\quad\hbox{i.e.}\quad y_0\le-0.50887956153832.
\]
Above that threshold there are two branches. Scaling both coordinates
by $H$, $X=x/H$ and $Y=u/H$ (so $dY/dX=u'$), gives
\[
\int_{-1}^1u\sqrt{1+u'^2}\,dx=H^2\int_{-1/H}^{1/H}Y\sqrt{1+Y'^2}\,dX,
\qquad Y(\pm1/H)=1:
\]
the catenoid problem of Exercise~22.1.6 with unit-ring half-separation
$1/H$. The scaled catenary $Y=(c/H)\cosh[X/(c/H)]$ has neck $c/H$, so
its endpoint parameter (half-separation over neck radius, called $u_0$
in Exercise~22.1.6) is $(1/H)/(c/H)=w$. Hence the smaller-$w$,
larger-$c$ branch ($w<u_c$) has positive second variation and is the
locally minimizing one.

\erratum{Part (b) has no single smooth minimizing ray for arbitrary
$y_0<1$. Also, Fermat's variational statement is generally stationarity;
an available stationary ray need not be the absolute minimum.}
The distinction can be quantified here. On the catenary
$u\sqrt{1+u'^2}=c\cosh^2(x/c)$, so, with $s=x/c$ and
$\int_0^w\cosh^2s\,ds=\tfrac12w+\tfrac14\sinh2w$, its optical length
is
\[
J_{\rm cat}=2ac^2\int_0^w\cosh^2s\,ds
=ac^2\left[w+\frac12\sinh(2w)\right],
\]
and since $c^2/H^2=\operatorname{sech}^2w$,
\[
\frac{J_{\rm cat}}{aH^2}
=\operatorname{sech}^2w\,(w+\sinh w\cosh w)
=w\operatorname{sech}^2w+\tanh w.
\]
To see the competing boundary cost, descend vertically from $H$
to a small $\delta>0$, travel the horizontal distance $2$ at that
height, and ascend to $H$. On the vertical segments $ds=|du|$. This
piecewise path has
\[
 J_\delta=a\left(2\int_\delta^H u\,du+2\delta\right)
 =a(H^2-\delta^2+2\delta)\longrightarrow aH^2.
\]
Smooth positive graphs approach the same limit. Let $u$ fall smoothly and
monotonically from $H$ to $\delta$ on $[-1,-1+\epsilon]$, equal
$\delta$ on $[-1+\epsilon,1-\epsilon]$, and rise likewise back to $H$
on $[1-\epsilon,1]$. Then $\int_{-1}^1u\,dx\le2\epsilon H+2\delta$,
since $u\le H$ on the two layers. Moreover $\sqrt{1+u'^2}\le1+|u'|$,
and on each monotone layer
$\int u|u'|\,dx=\int_\delta^Hu\,du=\tfrac12(H^2-\delta^2)$, so
\[
\frac Ja\le\int_{-1}^1u\,dx+\int_{-1}^1u|u'|\,dx
\le2\epsilon H+2\delta+H^2-\delta^2\longrightarrow H^2
\quad(\epsilon,\delta\to0).
\]
Equality of the catenary cost with this limit,
$w\operatorname{sech}^2w+\tanh w=1$, becomes after multiplication by
$2\cosh^2w$ the equation $2w+\sinh2w=2\cosh^2w$ of Exercise~22.1.7,
with the same root:
\[
2w-1=e^{-2w},\qquad w=0.639232271380537.
\]
On the shallow branch the ratio
$w\operatorname{sech}^2w+\tanh w$ has derivative
$2\operatorname{sech}^2w\,(1-w\tanh w)>0$, while
$H=\cosh w/w$ decreases as $w$ increases. So the ratio is below $1$
exactly when $w<0.639232\ldots$, i.e., when $H$ exceeds its value
$\cosh w/w=1/(w\operatorname{sech}w)=1/0.527697396962572$ at that root
(there $w\operatorname{sech}w$ is the half-separation $x_0$ of
Exercise~22.1.7): the shallow catenary beats that boundary limit
only for
\[
H>1.89502545541442.
\]
For a direct global proof, let $d=\min u>0$ on a competing
path. Cauchy--Schwarz applied to the vectors
$(d,\sqrt{u^2-d^2})$ and $(1,|u'|)$, of lengths $u$ and
$\sqrt{1+u'^2}$, gives
\[
 u\sqrt{1+u'^2}\ge d+\sqrt{u^2-d^2}\,|u'|.
\]
The path descends from $H$ to $d$ and returns to $H$. On a descending
or ascending stretch $|u'|\,dx=|du|$, so each of the two stretches
contributes at least $\int_d^H\sqrt{u^2-d^2}\,du$. With the
antiderivative
$\tfrac12u\sqrt{u^2-d^2}-\tfrac12d^2\operatorname{arcosh}(u/d)$,
integration over the horizontal span of length $2$ yields
\[
 \frac Ja\ge 2d+2\int_d^H\sqrt{u^2-d^2}\,du
 =:B(d)
 =2d+H\sqrt{H^2-d^2}-d^2\operatorname{arcosh}(H/d).
\]
For a path with additional rises and falls, the absolute-value
integral is only larger. Differentiation gives
\[
 B'(d)=2-2d\operatorname{arcosh}(H/d),\qquad B(0^+)=H^2:
\]
the $d$-derivative $-Hd/\sqrt{H^2-d^2}$ of $H\sqrt{H^2-d^2}$ cancels
the term $+Hd/\sqrt{H^2-d^2}$ that comes from differentiating
$\operatorname{arcosh}(H/d)$ in $-d^2\operatorname{arcosh}(H/d)$, and
$d^2\operatorname{arcosh}(H/d)\to0$ as $d\to0^+$ because
$\operatorname{arcosh}(H/d)\approx\ln(2H/d)$.

The product $d\operatorname{arcosh}(H/d)$ vanishes at both ends
$d=0,H$ and has a single maximum: putting
$d=H\operatorname{sech}w$ reduces it to $Hw\operatorname{sech}w$,
whose maximum is $0.662743\ldots H$. If $H\le1.50887956\ldots$, then
$B'\ge0$ on $(0,H)$, with at most one zero, so $B$ increases and
$B>H^2$ throughout. For larger $H$ there are two critical points:
$B'>0$ near $d=0$, $B'<0$ between them (where
$d\operatorname{arcosh}(H/d)>1$), and $B'>0$ beyond. So the smaller $d$
is a local maximum of $B$ and the larger $d$ is its local minimum. At
both critical points $d\operatorname{arcosh}(H/d)=1$, i.e.,
$H/d=\cosh(1/d)$: this is the catenary endpoint equation
$c\cosh(1/c)=H$ with $d=c$, and $w=\operatorname{arcosh}(H/d)=1/d$.
The two roots of $Hw\operatorname{sech}w=1$ lie on either side of $u_c$
(Exercise~22.1.6), and $d=H\operatorname{sech}w$ decreases as $w$
increases. So the larger root $d$, the minimum of $B$, has $w<u_c$: it
is the neck $c$ of the shallow catenary. Equality in the preceding bound
holds for that catenary, since there
$c+c\sinh^2(x/c)=c\cosh^2(x/c)=u\sqrt{1+u'^2}$. The endpoint $d=H$
cannot minimize $B$, since $B'(H^-)=2$. Thus, for
$H>1.50887956\ldots$, the global lower bound is
the smaller of the shallow-catenary cost and $H^2$. The former is at
most $H^2$, and is then attained by the catenary, precisely for
$H\ge1.8950254554144\ldots$; at equality it ties the boundary limit.
For smaller $H$, every $d>0$ gives a strictly larger cost than $H^2$,
so no smooth positive-index global minimizer exists. The infimum
$aH^2$ is approached by the endpoint-layer paths described above.

\erratum{In particular, the Instructor's Manual's claim of a global
catenary minimum at $y_0=-0.70$ is incorrect: $H=1.70$ lies below the
cost-crossing value $1.8950254554$.}


\Needspace{5\baselineskip}
\par\bigskip
\noindent\textbf{22.1.11}\quad A particle moves, starting at rest,
from point A on the surface of the Earth to point B (also on the
surface) by sliding frictionlessly through a tunnel. Find the
differential equation satisfied by the path if the transit time
is to be a minimum. Assume the Earth to be a nonrotating sphere
of uniform density.
\erratum{First take distinct nonantipodal endpoints, so $0<r_0<R$;
treat the antipodal diameter as a limit.}
\par\emph{Hint.}~The \erratum{gain in kinetic energy} of a particle of mass $m$ a
distance $r<R$ from the center of the Earth, with $R$ the Earth's
radius, is $\tfrac{1}{2}mg(R^{2}-r^{2})/R$, where $g$ is the
gravitational acceleration at the Earth's surface. It is
convenient to describe the path of the particle (in the plane
through A, B, and the center of the Earth) by plane polar
coordinates $(r,\theta)$, with A at $(R,-\varphi)$ and B at
$(R,\varphi)$.
\par\nobreak\hfill\emph{ANS.}\enspace
Letting $r_{0}$ be the minimum value of $r$ (reached at
$\theta=0$), Eq.~(22.21) yields
$r_{\theta}^{2}=\dfrac{r^{2}R^{2}(r^{2}-r_{0}^{2})}
{r_{0}^{2}(R^{2}-r^{2})}$ (the constant in that equation has the
value such that $r_{\theta}=0$ at $\theta=0$).
\par The solution for the path is a hypocycloid, generated by a
circle of radius $\tfrac{1}{2}(R-r_{0})$ rolling inside the
circle of radius $R$. You might like to show that the transit
time is
\[
t=\pi\frac{(R^{2}-r_{0}^{2})^{1/2}}{(Rg)^{1/2}}.
\]
For details see P.~W.~Cooper, \emph{Am.~J.~Phys.} \textbf{34}: 68
(1966); G.~Veneziano, et al., \textbf{34}: 701 (1966).
\solution
Inside a uniform sphere only the mass within radius $r$, the fraction
$r^3/R^3$ of the total, attracts the particle (Newton's shell theorem:
the uniform shell outside radius $r$ exerts no net force, and the inner
ball attracts as if its mass were at the center), so the surface force
$mg$ is multiplied by $(r^3/R^3)(R^2/r^2)=r/R$: the inward
gravitational force is $F_r=-mgr/R$. Thus
\[
U(r)-U(R)=-\int_R^rF_r\,dr'=\frac{mg}{R}\int_R^rr'\,dr'
=\frac{mg}{2R}(r^2-R^2).
\]
\erratum{The positive quantity in the printed hint is the
\emph{decrease} in potential energy from the surface, and hence the
kinetic energy gained. It is not the potential energy itself with
surface value zero: taken as $U(r)$, it would give the outward force
$-dU/dr=+mgr/R$.}
Energy conservation for release from rest gives
\[
\frac12mv^2=U(R)-U(r)=\frac{mg}{2R}(R^2-r^2),\qquad
v=\sqrt{\frac gR}\sqrt{R^2-r^2}.
\]
Writing the path as $r(\theta)$, the arc length is
$ds=\sqrt{dr^2+r^2d\theta^2}=\sqrt{r_\theta^2+r^2}\,d\theta$, and its
travel time $T=\int ds/v$ is
\[
T=\sqrt{\frac Rg}\int_{-\varphi}^{\varphi}
\frac{\sqrt{r_\theta^2+r^2}}{\sqrt{R^2-r^2}}\,d\theta.
\]
The integrand has no explicit $\theta$ dependence. With
$p=r_\theta$, omit the constant $\sqrt{R/g}$ and write
\[
f=\frac{\sqrt{r^2+p^2}}{\sqrt{R^2-r^2}},\qquad
f_p=\frac{p}{\sqrt{R^2-r^2}\sqrt{r^2+p^2}}.
\]
The Beltrami first integral
(Exercise~22.1.1 with $\theta$ in place of $x$) says that
\[
f-pf_p=\frac{(r^2+p^2)-p^2}
{\sqrt{R^2-r^2}\sqrt{r^2+p^2}}
=\frac{r^2}
{\sqrt{R^2-r^2}\sqrt{r^2+p^2}}
\]
is constant. At the minimum radius $r=r_0>0$, reached at $\theta=0$ by
symmetry, $r(\theta)$ is stationary, so $p=0$ there, and this
constant is $r_0^2/(r_0\sqrt{R^2-r_0^2})=r_0/\sqrt{R^2-r_0^2}$.
Squaring and rearranging,
\[
r^4(R^2-r_0^2)=r_0^2(R^2-r^2)(r^2+p^2),
\]
so that
\[
p^2=\frac{r^4(R^2-r_0^2)}{r_0^2(R^2-r^2)}-r^2
=\frac{r^2[r^2R^2-r^2r_0^2-r_0^2R^2+r_0^2r^2]}{r_0^2(R^2-r^2)},
\]
\[
\boxed{r_\theta^2=
\frac{r^2R^2(r^2-r_0^2)}{r_0^2(R^2-r^2)}}.
\]

To verify the stated geometric identification, put
$\rho=(R-r_0)/2$ and consider the hypocycloid
\[
z(u)=(R-\rho)e^{iu}+\rho\,e^{-i(R-\rho)u/\rho}.
\]
It is generated by a radius-$\rho$ circle rolling within the
radius-$R$ circle: its center $(R-\rho)e^{iu}$ turns through the angle $u$,
and rolling without slipping turns the tracing point through
$-(R-\rho)u/\rho$ relative to the center. Set $\delta=Ru/\rho$, and note
$u+(R-\rho)u/\rho=\delta$. Then
\[
\begin{aligned}
r^2=|z|^2&=(R-\rho)^2+\rho^2+2\rho(R-\rho)\cos\delta\\
&=R^2-2\rho(R-\rho)(1-\cos\delta)=R^2-4\rho(R-\rho)\sin^2(\delta/2),
\end{aligned}
\]
and differentiating, $2r\,dr/du=-2\rho(R-\rho)\sin\delta\,(R/\rho)$. For the
polar angle $\theta=\arg z$, write $z=re^{i\theta}$; then
$z'=dz/du=(r'+ir\theta')e^{i\theta}$ and $\bar z\,z'=rr'+ir^2\theta'$, so
$d\theta/du=\operatorname{Im}(\bar z\,z')/|z|^2$. Here
$z'=i(R-\rho)(e^{iu}-e^{-i(R-\rho)u/\rho})$. Multiplying out,
\[
\bar z\,z'=i(R-\rho)\left[(R-\rho)(1-e^{-i\delta})+\rho(e^{i\delta}-1)\right],
\]
so $\operatorname{Im}(\bar z\,z')=(R-\rho)(R-2\rho)(1-\cos\delta)$, and
$R-2\rho=r_0$. Hence
\[
\frac{dr}{du}=-\frac{R(R-\rho)\sin\delta}{r},\qquad
\frac{d\theta}{du}
=\frac{(R-\rho)r_0(1-\cos\delta)}{r^2}.
\]
Their ratio is
$dr/d\theta=-(Rr/r_0)\sin\delta/(1-\cos\delta)=-(Rr/r_0)\cot(\delta/2)$.
Since $4\rho(R-\rho)=(R-r_0)(R+r_0)=R^2-r_0^2$, the expression for $r^2$
gives $\sin^2(\delta/2)=(R^2-r^2)/(R^2-r_0^2)$ and
$\cos^2(\delta/2)=(r^2-r_0^2)/(R^2-r_0^2)$. Therefore
\[
\left(\frac{dr}{d\theta}\right)^2
=\frac{R^2r^2}{r_0^2}\cot^2(\delta/2)
=\frac{R^2r^2(r^2-r_0^2)}
{r_0^2(R^2-r^2)},
\]
the boxed equation. The radius is $r_0$ at $\delta=\pi$ and $R$ at the
cusps $\delta=0,2\pi$. Rotating this arc puts its minimum
radius at $\theta=0$ and its two surface cusps at $\pm\varphi$.

The transit time can be evaluated without integrating the path explicitly.
On either half the boxed equation gives $|d\theta/dr|$, and
$ds/dr=\sqrt{1+r^2(d\theta/dr)^2}$, where
$R^2(r^2-r_0^2)+r_0^2(R^2-r^2)=r^2(R^2-r_0^2)$:
\[
\left|\frac{d\theta}{dr}\right|=
\frac{r_0\sqrt{R^2-r^2}}{Rr\sqrt{r^2-r_0^2}},\qquad
\frac{ds}{dr}=
\frac{r\sqrt{R^2-r_0^2}}{R\sqrt{r^2-r_0^2}}.
\]
Dividing $ds$ by $v=\sqrt{g/R}\sqrt{R^2-r^2}$ and doubling for the
two halves, with $\Delta=R^2-r_0^2$,
\[
T=2\sqrt{\frac Rg}\frac{\sqrt\Delta}{R}
\int_{r_0}^R\frac{r\,dr}
{\sqrt{r^2-r_0^2}\sqrt{R^2-r^2}}.
\]
Set $r^2=r_0^2+\Delta\sin^2\psi$. Then $r\,dr=\Delta\sin\psi\cos\psi\,d\psi$,
$r^2-r_0^2=\Delta\sin^2\psi$ and $R^2-r^2=\Delta\cos^2\psi$, so the
integrand becomes $d\psi$, and the limits $r_0,R$ become $0,\pi/2$.
Hence $T=2\sqrt{R/g}\,(\sqrt\Delta/R)(\pi/2)$, that is,
\[
\boxed{T=\pi\sqrt{\frac{R^2-r_0^2}{Rg}}}.
\]
The same substitution in the angular integral
$\varphi=\int_{r_0}^R|d\theta/dr|\,dr$ turns the integrand into
$(r_0/R)\Delta\cos^2\psi\,d\psi/r^2$, where
$r^2=r_0^2\cos^2\psi+R^2\sin^2\psi$ and
$\Delta\cos^2\psi=R^2-r^2$. Hence
\[
\varphi=\frac{r_0}{R}\left[
R^2\int_0^{\pi/2}\frac{d\psi}
{r_0^2\cos^2\psi+R^2\sin^2\psi}-\frac\pi2\right]
=\frac\pi2\left(1-\frac{r_0}{R}\right).
\]
Here the substitution $t=\tan\psi$ turns the integral into
$\int_0^\infty dt/(r_0^2+R^2t^2)=\pi/(2r_0R)$. As a check, one arch
of the hypocycloid spans the arc $2\pi\rho$ of the large circle, i.e.,
the angle $2\varphi=2\pi\rho/R=\pi(1-r_0/R)$.
The antipodal case is the limiting diameter $r_0=0$,
with $T=\pi\sqrt{R/g}$; the displayed polar ODE assumes $r_0>0$.


\Needspace{5\baselineskip}
\par\bigskip
\noindent\textbf{22.1.12}\quad A ray of light follows a
straight-line path in a first homogeneous medium, is refracted
at an interface, and then follows a new straight-line path in the
second medium. See Fig.~22.7. Use Fermat's principle of optics
to derive Snell's law of refraction:
\[
n_{1}\sin\theta_{1}=n_{2}\sin\theta_{2}.
\]
\par\emph{Hint.}~Keep the points $(x_{1},y_{1})$ and
$(x_{2},y_{2})$ fixed and vary $x_{0}$ to satisfy Fermat's
principle.
\par\emph{Note.}~This is not an Euler equation problem, because
the light path is not differentiable at $x_{0}$.
\solution
Take the interface to be $y=0$, with fixed endpoints
$(x_1,y_1)$ above it and $(x_2,y_2)$ below it. The crossing point
is $(x_0,0)$. In a medium of index $n$ light travels at $c/n$, with
$c$ the vacuum light speed, so each straight segment of length $L$ takes
the time $nL/c$, and the travel time is $J(x_0)/c$ with the optical
path length
\[
J(x_0)=n_1\sqrt{(x_0-x_1)^2+y_1^2}
+n_2\sqrt{(x_2-x_0)^2+y_2^2}.
\]
Differentiate only the crossing position:
\[
J'(x_0)=
n_1\frac{x_0-x_1}{\sqrt{(x_0-x_1)^2+y_1^2}}
-n_2\frac{x_2-x_0}{\sqrt{(x_2-x_0)^2+y_2^2}}.
\]
The interface normal is the $y$ direction. The incident segment has
horizontal extent $x_0-x_1$ and length $\sqrt{(x_0-x_1)^2+y_1^2}$, so
the first ratio is $\sin\theta_1$, with $\theta_1$ the angle between
the incident ray and the normal; likewise the second ratio is
$\sin\theta_2$ for the refracted ray (both sines signed, positive when
the ray advances toward $+x$). Thus $J'(x_0)=0$ gives
\[
\boxed{n_1\sin\theta_1=n_2\sin\theta_2}.
\]
For positive indices and endpoints off the interface,
differentiating $s/\sqrt{s^2+y^2}$ gives $y^2/(s^2+y^2)^{3/2}$, so
\[
J''(x_0)=
\frac{n_1y_1^2}{[(x_0-x_1)^2+y_1^2]^{3/2}}
+\frac{n_2y_2^2}{[(x_2-x_0)^2+y_2^2]^{3/2}}>0.
\]
Hence $J'$ increases strictly from $-(n_1+n_2)$ as $x_0\to-\infty$ to
$n_1+n_2$ as $x_0\to+\infty$; it has exactly one zero, where the
strictly convex $J$ has its minimum. This proves the stationary
crossing is the unique minimum for these fixed endpoints. No
differentiability across the corner was assumed.


\Needspace{5\baselineskip}
\par\bigskip
\noindent\textbf{22.1.13}\quad A second soap-film configuration
for the unit-radius rings at $x=\pm x_{0}$ consists of a circular
disk, radius $a$, in the $x=0$ plane and two catenoids of
revolution, one joining the disk and each ring. One catenoid may
be described by
\[
y=c_{1}\cosh\left(\frac{x}{c_{1}}+c_{3}\right).
\]
\begin{enumerate}
\item[(a)] Impose boundary conditions at $x=0$ and $x=x_{0}$.
\item[(b)] Require that the catenoids form an angle of $120^{\circ}$
where they join the central disk. Express this third boundary
condition mathematically. \erratum{For films of equal surface tension
and a freely moving junction, this is the equilibrium condition.}
\item[(c)] Show that the total area of catenoids plus central
disk is then
\[
A=\erratum{\pi}c_{1}^{2}\left[\sinh\left(\frac{2x_{0}}{c_{1}}+2c_{3}\right)
+\frac{2x_{0}}{c_{1}}\right].
\]
\end{enumerate}
\par\emph{Note.}~\erratum{Stability depends on the branch and the
allowed variations; see Exercise~22.1.14. The area is larger than
that of the \emph{shallow} simple catenoid whenever both exist.}
\par\nobreak\hfill\emph{ANS.}\enspace
(a) $\begin{cases}
1=c_{1}\cosh\left(\dfrac{x_{0}}{c_{1}}+c_{3}\right)\\[8pt]
a=c_{1}\cosh c_{3}
\end{cases}$
\quad(b) $\dfrac{dy}{dx}=\tan30^{\circ}=\sinh c_{3}$.
\solution
(a) Write $c=c_1$. The right-hand catenoid
$y=c\cosh(x/c+c_3)$, $0\le x\le x_0$, has radius $a$ where it meets
the disk ($x=0$) and radius $1$ at the ring ($x=x_0$), so the boundary
equations are
\[
\boxed{a=c\cosh c_3,\qquad
1=c\cosh(x_0/c+c_3)}.
\]

(b) Work in a meridian plane through the axis, where the three films
appear as curves meeting at $(0,a)$. At the junction, the central disk
points inward in the negative radial direction, with unit tangent
$(0,-1)$. The outward right-catenoid unit tangent
is $(1,y')/\sqrt{1+y'^2}$. A $120^\circ$ angle therefore requires
\[
-\frac{y'}{\sqrt{1+y'^2}}=\cos120^\circ=-\frac12,
\]
that is, $4y'^2=1+y'^2$, which gives $y'=1/\sqrt3=\tan30^\circ$.
By reflection the left catenoid has unit tangent
$(-1,y')/\sqrt{1+y'^2}$; its angle with the right one has cosine
$(y'^2-1)/(y'^2+1)=-1/2$, so all three angles are then $120^\circ$.
Since $y'=\sinh(x/c+c_3)$, at $x=0$ this is $\sinh c_3=1/\sqrt3$.
Then $\cosh c_3=\sqrt{1+1/3}=2/\sqrt3$, and
$\operatorname{arsinh}s=\ln(s+\sqrt{1+s^2})$ gives
$c_3=\ln(1/\sqrt3+2/\sqrt3)=\ln\sqrt3$:
\[
\boxed{\sinh c_3=\frac1{\sqrt3},\qquad
c_3=\operatorname{arsinh}(1/\sqrt3)=\tfrac12\ln3,\qquad
a=c\cosh c_3=\frac{2c}{\sqrt3}}.
\]

The same angle follows by varying the junction radius. In the
reflection-symmetric class the area is
$A=4\pi\int_0^{x_0}y\sqrt{1+y'^2}\,dx+\pi y(0)^2$.
Vary $y\to y+\alpha\eta$ with $\eta(x_0)=0$ and $\eta(0)$ free, so
that the junction radius changes by $\delta a=\alpha\eta(0)$. For
$f=y\sqrt{1+y'^2}$, $f_{y'}=yy'/\sqrt{1+y'^2}$, and the integration by
parts of Exercise~22.1.2 gives, to first order in $\alpha$,
\[
\delta\int_0^{x_0}f\,dx=\alpha\bigl[f_{y'}\eta\bigr]_0^{x_0}
+\alpha\int_0^{x_0}\left(f_y-\frac{d}{dx}f_{y'}\right)\eta\,dx.
\]
The integral vanishes by the interior Euler equation, the upper
boundary term vanishes because $\eta(x_0)=0$, and the lower one is
$-\alpha f_{y'}(0)\eta(0)=-a\,y'(0)\,\delta a/\sqrt{1+y'(0)^2}$,
which enters $A$ with the factor $4\pi$. The disk contributes
$\delta(\pi a^2)=2\pi a\,\delta a$. So the first variation is
\[
 \delta A=\left[-4\pi a\frac{y'(0)}{\sqrt{1+y'(0)^2}}
                    +2\pi a\right]\delta a.
\]
Since $a>0$ and $\delta a$ is free, this again gives
$y'(0)/\sqrt{1+y'(0)^2}=1/2$. Thus the angle is necessary for
this free-junction equilibrium, although one can join mathematical
surface pieces at other prescribed angles.

(c) Put $v=x_0/c+c_3$. There are two catenoids and one disk. On each
catenoid $\sqrt{1+y'^2}=\cosh(x/c+c_3)$, so
$y\sqrt{1+y'^2}=c\cosh^2(x/c+c_3)$. Substitute $u=x/c+c_3$,
$dx=c\,du$ (limits $c_3$ and $v$), and use
$\int\cosh^2u\,du=\tfrac12u+\tfrac14\sinh2u$ and $a=c\cosh c_3$.
Their total area is
\[
\begin{aligned}
A&=4\pi\int_0^{x_0}y\sqrt{1+y'^2}\,dx+\pi a^2\\
&=4\pi c^2\int_{c_3}^{v}\cosh^2u\,du+\pi c^2\cosh^2c_3\\
&=\pi c^2[2(v-c_3)+\sinh2v-\sinh2c_3+\cosh^2c_3].
\end{aligned}
\]
The angle condition gives
$\cosh c_3=2/\sqrt3$, so
$\sinh2c_3=2\sinh c_3\cosh c_3=2\cdot\frac1{\sqrt3}\cdot\frac2{\sqrt3}
=\frac43=\cosh^2c_3$. These last terms cancel, and
$2(v-c_3)=2x_0/c$, $2v=2x_0/c+2c_3$:
\[
\boxed{A=\pi c^2\left[
\sinh\left(\frac{2x_0}{c}+2c_3\right)+\frac{2x_0}{c}
\right]}.
\]
\erratum{The printed expression is missing the overall factor $\pi$.
It equals $A/\pi$, not the total physical area. The surface-of-revolution
factor $2\pi y\,ds$ and the disk area $\pi a^2$ fix the normalization.}


\Needspace{5\baselineskip}
\par\bigskip
\noindent\textbf{22.1.14}\quad For the soap film described in
Exercise~22.1.13, find (numerically) the maximum value of
$x_{0}$.
\par\emph{Note.}~This calls for a calculator with hyperbolic
functions or a table of hyperbolic cotangents.
\par\nobreak\hfill\emph{ANS.}\enspace$x_{0\max}=0.4078$.
\solution
Retain the angle condition from Exercise~22.1.13:
$c_3=\tfrac12\ln3$. Write $v=x_0/c+c_3$. The ring condition
$1=c\cosh v$ of Exercise~22.1.13(a) gives $c=\operatorname{sech}v$, and
then $x_0=c(v-c_3)$:
\[
v=\frac{x_0}{c}+c_3,\qquad
c=\operatorname{sech}v,\qquad
x_0=(v-c_3)\operatorname{sech}v,\quad v>c_3.
\]
Differentiate by the product rule, using
$(\operatorname{sech}v)'=-\operatorname{sech}v\tanh v$:
\[
\frac{dx_0}{dv}
=\operatorname{sech}v[1-(v-c_3)\tanh v].
\]
Thus the maximum satisfies
\[
\boxed{v-c_3=\coth v}.
\]
The function $v-c_3-\coth v$ has derivative
$1+\operatorname{csch}^2v>0$, equals $-\coth c_3=-2$ at $v=c_3$
(because $\tanh c_3=\sinh c_3/\cosh c_3=1/2$),
and tends to infinity. It has exactly one root. Since
$1-(v-c_3)\tanh v=\tanh v\,[\coth v-(v-c_3)]$, $x_0$ increases before
the root and decreases after it. Moreover $x_0$ tends to zero at both
ends of the parameter range ($x_0=0$ at $v=c_3$, and
$x_0<v\operatorname{sech}v\to0$ as $v\to\infty$), so this root is the
global maximum.

Newton's update is
\[
v_{j+1}=v_j-\frac{v_j-c_3-\coth v_j}
{1+\operatorname{csch}^2v_j}.
\]
Starting with $v_0=1.6$, three steps give all the digits shown:
\[
v=1.62926936306647,\qquad
c=\operatorname{sech}v=0.377627833193942,
\]
and hence
\[
\boxed{x_{0\max}=(v-c_3)\operatorname{sech}v
=0.407824170219077}.
\]
This rounds to the quoted answer $0.4078$.

\erratum{The two branches must not both be called stable.}
To check this, allow reflection-symmetric, axisymmetric variations
of the two catenoids, with the junction remaining in $x=0$ and its
radius free. Set $u=x/c+c_3$ and $p(u)=\operatorname{sech}^2u$.
For $y\to y+\alpha\eta$ on the right catenoid (the left one is its
mirror image), the second variation of $4\pi\int_0^{x_0}y\sqrt{1+y'^2}\,dx$
is, with the second derivatives of $f=y\sqrt{1+y'^2}$ from
Exercise~22.1.6 (now $f_{yy'}=\tanh u$, $f_{y'y'}=c\operatorname{sech}^2u$),
\[
4\pi\int_0^{x_0}\left[2\tanh u\,\eta\eta_x
+c\operatorname{sech}^2u\,\eta_x^2\right]dx.
\]
With $dx=c\,du$ and $\eta_x=\eta_u/c$, the mixed term becomes
$4\pi\int_{c_3}^v\tanh u\,(\eta^2)_u\,du$ and the last term
$4\pi\int_{c_3}^vp\,\eta_u^2\,du$. Integrating the former by parts,
with $(\tanh u)'=p$ and $\eta(v)=0$, the catenoids contribute
\[
4\pi\left[\int_{c_3}^vp\,(\eta_u^2-\eta^2)\,du
-\tanh(c_3)\,\eta(c_3)^2\right],\qquad \tanh c_3=\tfrac12.
\]
The disk area $\pi(a+\alpha\eta(c_3))^2$ has second $\alpha$-derivative
$2\pi\eta(c_3)^2=4\pi\cdot\tfrac12\eta(c_3)^2$. The boundary terms
cancel, so
\[
 Q[\eta]=\frac{\delta^2A}{4\pi}
 =\int_{c_3}^{v}p(u)(\eta'^2-\eta^2)\,du,
 \qquad \eta(v)=0,\quad \eta(c_3)\text{ free}.
\]
The function
\[
 w(u)=\cosh u-(u-c_3)\sinh u,
\]
which is $\partial y/\partial c$ for the family $y=c\cosh(x/c+c_3)$ at
fixed $x$ and $c_3$, has $w'=-(u-c_3)\cosh u$. Hence $w'(c_3)=0$,
$pw'=-(u-c_3)\operatorname{sech}u$, and
$(pw')'=-\operatorname{sech}u+(u-c_3)\operatorname{sech}u\tanh u=-pw$,
so $(pw')'+pw=0$. Since $w=\sinh u\,[\coth u-(u-c_3)]$ and the bracket
decreases strictly from $\coth c_3=2$, the first zero of $w$ is
$v_c=1.629269\ldots$, the root just found. For $v<v_c$, $w>0$ on
$[c_3,v]$; putting $\eta=wz$ gives the same square identity as in
Exercise~22.1.6,
$p(\eta'^2-\eta^2)=pw^2z'^2+(pww'z^2)'$. The boundary term
$[pww'z^2]_{c_3}^v$ vanishes at $u=v$ because $z(v)=\eta(v)/w(v)=0$,
and at $u=c_3$ because $w'(c_3)=0$. So
\[
Q[\eta]=\int_{c_3}^{v}pw^2z'^2\,du>0\qquad(\eta\ne0),
\]
since $z'\equiv0$ and $z(v)=0$ would force $\eta\equiv0$.
To pass from this quadratic test to a local minimum, recall that the
first variation vanishes for these variations, by the interior Euler
equation and the $120^\circ$ junction condition (Exercise~22.1.13(b)).
So $A[y+\eta]-A[y]=\tfrac12\delta^2A+R=2\pi Q[\eta]+R$, with $R$ the
Taylor remainder. Next, note that $z(v)=0$ implies
$z(u)=-\int_u^v z'(s)ds$, so
$z(u)^2\le(v-c_3)\int_{c_3}^vz'^2ds$ by Cauchy--Schwarz, and
$\int z^2\le(v-c_3)^2\int z'^2$.
On $[c_3,v]$ put $m=\min pw^2>0$ and $M=\max(w^2+2w'^2,2w^2)$. Since
$\eta^2+\eta'^2=w^2z^2+(w'z+wz')^2\le M(z^2+z'^2)$,
\[
\int(\eta^2+\eta'^2)du\le M[1+(v-c_3)^2]\int z'^2du
\le\frac{M[1+(v-c_3)^2]}{m}\,Q[\eta],
\]
that is, $Q[\eta]\ge\kappa\int(\eta^2+\eta'^2)du$
with $\kappa=m/\{M[1+(v-c_3)^2]\}>0$. For $\eta\ne0$ with sufficiently
small $\|\eta\|_{C^1}$, the remainder satisfies
$|R|<2\pi\kappa\int(\eta^2+\eta'^2)du\le2\pi Q[\eta]$, by uniform
continuity of the integrand's Hessian (the disk term is exactly
quadratic), so $A[y+\eta]>A[y]$. This proves strict local minimality on
the shallow branch in the stated class. At $v=v_c$, $\eta=w$ is a
zero mode; the second variation alone does not classify higher orders.

For $v>v_c$, extend $w$ by zero from $v_c$ to $v$, calling the
result $\eta_0$. It is an admissible piecewise smooth variation
($\eta_0(v)=0$), and integration by parts on $[c_3,v_c]$ gives
$Q[\eta_0]=[pww']_{c_3}^{v_c}-\int_{c_3}^{v_c}w[(pw')'+pw]\,du=0$,
because $w(v_c)=0=w'(c_3)$. Choose a smooth $h$ with $h(v)=0$ and
$h(v_c)>0$. Integration by parts on $[c_3,v_c]$ gives
\[
 Q[\eta_0,h]
 :=\int_{c_3}^{v}p(\eta_0'h'-\eta_0h)\,du
 =\bigl[pw'h\bigr]_{c_3}^{v_c}-\int_{c_3}^{v_c}[(pw')'+pw]h\,du
 =p(v_c)w'(v_c)h(v_c)<0,
\]
since $w'(c_3)=0$ and $w'(v_c)=-(v_c-c_3)\cosh v_c<0$.
Therefore $Q[\eta_0+\epsilon h]=Q[\eta_0]
+2\epsilon Q[\eta_0,h]+\epsilon^2Q[h]<0$ for small positive
$\epsilon$. Smoothing the corner preserves this strict inequality,
because $Q$ changes little when $\eta$ and $\eta'$ change little in
mean square. The deep branch is unstable even in this restricted class.

The area comparison in Exercise~22.1.13 also needs its branch specified.
Write the simple catenoid of Exercise~22.1.6 as
$x_0=\mu\operatorname{sech}\mu$, where $\mu$ is the parameter called
$u$ there: the ratio of $x_0$ to that catenoid's own neck radius
$\operatorname{sech}\mu$, not to the three-film $c$ above. By
Exercise~22.1.7 its area is
$A_s=\pi\operatorname{sech}^2\mu\,(2\sinh\mu\cosh\mu+2\mu)
=2\pi(\tanh\mu+\mu\operatorname{sech}^2\mu)$. Take its shallow
root $\mu<u_c$ from Exercise~22.1.6. Let $A_t$ denote the three-film
area; Exercise~22.1.13(c) with $c=\operatorname{sech}v$ and
$2x_0/c=2(v-c_3)$ gives
$A_t=2\pi[\tanh v+(v-c_3)\operatorname{sech}^2v]$. Differentiating,
\[
\frac{dA_s}{d\mu}=4\pi\operatorname{sech}^2\mu\,(1-\mu\tanh\mu),\qquad
\frac{dA_t}{dv}=4\pi\operatorname{sech}^2v\,[1-(v-c_3)\tanh v],
\]
and dividing by $dx_0/d\mu$ (Exercise~22.1.6(b)) and by $dx_0/dv$
(above) gives, along either family,
\[
 \frac{dA_t}{dx_0}=4\pi\operatorname{sech}v,\qquad
 \frac{dA_s}{dx_0}=4\pi\operatorname{sech}\mu.
\]
At a common separation, $\mu<v$: for $v\le u_c$, use the increasing
function $\mu\operatorname{sech}\mu$ on $(0,u_c]$ and
$\mu\operatorname{sech}\mu=x_0=(v-c_3)\operatorname{sech}v
<v\operatorname{sech}v$;
for $v>u_c$ it is immediate, since $\mu<u_c$. As $\operatorname{sech}$
decreases, $d(A_t-A_s)/dx_0=4\pi(\operatorname{sech}v
-\operatorname{sech}\mu)<0$: $A_t-A_s$ decreases with $x_0$
on the shallow three-film branch. At its fold, where
$v_c-c_3=\coth v_c$ and
$\tanh v+\coth v\operatorname{sech}^2v=(\sinh^2v+1)/(\sinh v\cosh v)
=\coth v$,
\[
 A_t=2\pi\bigl[\tanh v_c+\coth v_c\operatorname{sech}^2v_c\bigr]
     =2\pi\coth v_c>2\pi.
\]
Meanwhile $x_{0\max}<1/2<0.527697\ldots$, the area crossing in
Exercise~22.1.7; since $A_s$ increases with $x_0$ on the shallow branch
and equals $2\pi$ at the crossing, $A_s<2\pi$ throughout this range.
The bound $x_0<1/2$, i.e. $\tfrac12\cosh v-v+c_3>0$, follows without
rounding from $\cosh v\ge1+v^2/2$ and $c_3=\tfrac12\ln3>\tfrac12$
(as $3>e$):
\[
 \tfrac12\cosh v-v+c_3
 \ge\tfrac12+\tfrac14v^2-v+c_3
 =\tfrac14(v-2)^2+c_3-\tfrac12>0.
\]
The shallow-branch difference is therefore positive even at its
minimum, the fold. On the deep three-film branch,
$v-c_3>\coth v$, so $A_t>2\pi\coth v>2\pi>A_s$ directly.
This proves the stated comparison for both three-film branches.
It is not a comparison with the deep simple catenoid.


\Needspace{5\baselineskip}
\par\bigskip
\noindent\textbf{22.1.15}\quad Find the curve of quickest descent
from $(0,0)$ to $(x_{0},y_{0})$ for a particle that, starting
from rest, slides under gravity and without friction.
\erratum{Assume $x_0,y_0>0$ with $y$ measured downward.} \erratum{For
the special endpoint ratio $x_0/y_0=\pi/2$,} show that
the ratio of times taken by the particle along a straight line
joining the two points compared to along the curve of quickest
descent is $(1+4/\pi^{2})^{1/2}$.
\par\emph{Hint.}~Take $y$ to increase downwards. Apply
Eq.~(22.21) to obtain $y_{x}^{2}=(1-c^{2}y)/c^{2}y$, where $c$ is
an integration constant. It is helpful to make the substitution
$c^{2}y=\sin^{2}(\varphi/2)$ and take
$(x_{0},y_{0})=(\pi/2c^{2},\,1/c^{2})$.
\solution
Take $y$ positive downward. Conservation of energy,
$\tfrac12mv^2=mgy$ for a start from rest at $y=0$, gives
$v=\sqrt{2gy}$, so $T=\int ds/v$ with $ds=\sqrt{1+y'^2}\,dx$ is
\[
T[y]=\frac1{\sqrt{2g}}\int
\frac{\sqrt{1+y'^2}}{\sqrt y}\,dx.
\]
For $f=\sqrt{1+y'^2}/\sqrt y$, $f_{y'}=y'/(\sqrt y\sqrt{1+y'^2})$ and
$f-y'f_{y'}=[(1+y'^2)-y'^2]/(\sqrt y\sqrt{1+y'^2})$. Since $f$ does not
depend on $x$, this is constant (Eq.~(22.21); Exercise~22.1.1 with
$f_x=0$). The Beltrami first integral, with its constant written
$c>0$, is therefore
\[
\frac1{\sqrt y\sqrt{1+y'^2}}=c,\qquad
y'^2=\frac{1-c^2y}{c^2y},
\]
the second form following from $y(1+y'^2)=1/c^2$.
Put $c^2y=\sin^2(\varphi/2)$. Then $1-c^2y=\cos^2(\varphi/2)$, so
$y'^2=\cot^2(\varphi/2)$ and $y'=\cot(\varphi/2)$: positive while the
particle descends ($0<\varphi<\pi$) and negative beyond the lowest
point. Also
\[
y=\frac{1-\cos\varphi}{2c^2},\qquad
\frac{dy}{d\varphi}=\frac{\sin\varphi}{2c^2},\qquad
\frac{dx}{d\varphi}=\frac{dy/d\varphi}{y'}
=\frac{\sin\varphi\tan(\varphi/2)}{2c^2}
=\frac{1-\cos\varphi}{2c^2},
\]
using $\sin\varphi\tan(\varphi/2)=2\sin^2(\varphi/2)=1-\cos\varphi$.
The start $(0,0)$ has $y=0$, i.e. $\varphi=0$. Choosing $x=0$ at
$\varphi=0$ and integrating gives the cycloid
\[
\boxed{x=\frac{\varphi-\sin\varphi}{2c^2},\qquad
y=\frac{1-\cos\varphi}{2c^2}}.
\]
For arbitrary $x_0,y_0>0$, choose $0<\varphi_0<2\pi$ from
\[
\frac{x_0}{y_0}
=\frac{\varphi_0-\sin\varphi_0}{1-\cos\varphi_0},
\qquad
\frac1{2c^2}=\frac{y_0}{1-\cos\varphi_0}.
\]
The ratio increases from zero to infinity on this interval. Near
$\varphi=0$ it behaves like $(\varphi^3/6)/(\varphi^2/2)=\varphi/3$, and
as $\varphi\to2\pi$ its denominator tends to zero while its numerator
tends to $2\pi$. By the quotient rule, with
$(1-\cos\varphi)^2+\sin^2\varphi=2(1-\cos\varphi)$, its derivative
is $[2(1-\cos\varphi)-\varphi\sin\varphi]/(1-\cos\varphi)^2$, and the
numerator is positive. For $0<\varphi<\pi$ it equals
\[
4\sin^2\frac\varphi2-2\varphi\sin\frac\varphi2\cos\frac\varphi2
=2\sin\frac\varphi2\cos\frac\varphi2
\left(2\tan\frac\varphi2-\varphi\right)>0,
\]
because $\tan t>t$ for $0<t<\pi/2$; for $\pi\le\varphi<2\pi$,
$\sin\varphi\le0$ and the sign is immediate. Hence $\varphi_0$ is
unique, and then $c$ follows from the second equation.

With $A=1/(2c^2)$, $x=A(\varphi-\sin\varphi)$ and
$y=A(1-\cos\varphi)$, so
$(dx/d\varphi)^2+(dy/d\varphi)^2
=A^2[(1-\cos\varphi)^2+\sin^2\varphi]=4A^2\sin^2(\varphi/2)$
and $2gy=4gA\sin^2(\varphi/2)$:
\[
ds=2A\sin(\varphi/2)d\varphi,\qquad
v=2\sqrt{gA}\sin(\varphi/2),
\]
so $dt=ds/v=\sqrt{A/g}\,d\varphi$ and
\[
T_{\rm cyc}=\varphi_0\sqrt{A/g}.
\]
One can verify the minimizing character by a calibration.
Regard $S(x,y)=\varphi\sqrt{A/g}$ as a function defined by the
first-arch endpoint equations $x=A(\varphi-\sin\varphi)$,
$y=A(1-\cos\varphi)$, $A>0$, $0<\varphi<2\pi$. By the monotonicity
just shown, each point with $x,y>0$ has exactly one such
$(A,\varphi)$, and the Jacobian
$\partial(x,y)/\partial(A,\varphi)
=-A[2(1-\cos\varphi)-\varphi\sin\varphi]$ is nonzero, so $S$ is smooth
there. The differentials of the endpoint equations are
\[
 dx=(\varphi-\sin\varphi)\,dA+A(1-\cos\varphi)\,d\varphi,\qquad
 dy=(1-\cos\varphi)\,dA+A\sin\varphi\,d\varphi,
\]
whereas
\[
dS=\frac{\varphi\,dA}{2\sqrt{gA}}+\sqrt{\frac Ag}\,d\varphi.
\]
Equating the coefficients of $dA,d\varphi$ in $dS=S_xdx+S_ydy$ gives
two linear equations for $S_x,S_y$, solved by
\[
S_x=\frac1{2\sqrt{gA}},\qquad
S_y=\frac{\cot(\varphi/2)}{2\sqrt{gA}}.
\]
Indeed, use $(1-\cos\varphi)\cot(\varphi/2)=\sin\varphi$ in the first
coefficient and $\sin\varphi\cot(\varphi/2)=1+\cos\varphi$ in the
second. Consequently, since $2gy=4gA\sin^2(\varphi/2)$,
\[
|\nabla S|^2=\frac{1+\cot^2(\varphi/2)}{4gA}
=\frac{\csc^2(\varphi/2)}{4gA}=\frac1{2gy}.
\]
Also $S\to0$ as $(x,y)\to(0,0)$, by the following two bounds. When
$\varphi\le\pi$, eliminating $A$ with $y=2A\sin^2(\varphi/2)$ gives
\[
S=\frac{\varphi}{\sin(\varphi/2)}\sqrt{\frac y{2g}}
\le\pi\sqrt{\frac y{2g}},
\]
because $\sin(\varphi/2)\ge\varphi/\pi$: the sine is concave on
$[0,\pi/2]$, so there its graph lies above the chord $2t/\pi$ from
$(0,0)$ to $(\pi/2,1)$; take $t=\varphi/2$. When $\pi\le\varphi<2\pi$, $\sin\varphi\le0$ gives
$x=A(\varphi-\sin\varphi)\ge\pi A$, so
$S=\varphi\sqrt{A/g}<2\pi\sqrt{x/(\pi g)}=2\sqrt{\pi x/g}$.

For any joining path in $y>0$, Cauchy--Schwarz gives
$dS=\nabla S\cdot d\mathbf r\le|\nabla S|\,ds=ds/\sqrt{2gy}=dt$;
integrating from the initial limiting point, where $S=0$, proves
$T\ge S(x_0,y_0)=T_{\rm cyc}$. Equality holds along the cycloid, whose
tangent $(dx,dy)\propto(1-\cos\varphi,\sin\varphi)
\propto(1,\cot(\varphi/2))$ is parallel to $\nabla S$.
Reflection handles paths to negative $x$: replacing $x$ by $|x|$
changes neither the arc length nor the speed.

\erratum{The fixed numerical time ratio in the question applies
to the special endpoint in the hint, not to arbitrary $(x_0,y_0)$.}
For that endpoint, $x_0/y_0=\pi/2$ is the value of the ratio at
$\varphi_0=\pi$, so $\varphi_0=\pi$ and $A=y_0/2$: indeed
$x_0=\pi/(2c^2)$, $y_0=1/c^2$, and
$T_{\rm cyc}=\pi\sqrt{A/g}=\pi/(c\sqrt{2g})$.
The straight segment has length
\[
L=\sqrt{x_0^2+y_0^2}
=\frac1{c^2}\sqrt{1+\pi^2/4}
\]
and constant tangential acceleration $g y_0/L$ ($g$ times the sine of
its inclination). Starting from rest, $L=\tfrac12(gy_0/L)T_{\rm line}^2$,
so, with $\sqrt{2/(gy_0)}=c\sqrt{2/g}$,
\[
T_{\rm line}=\sqrt{\frac{2L^2}{g y_0}}
=\frac1c\sqrt{\frac2g}\sqrt{1+\frac{\pi^2}{4}}.
\]
Therefore, since $\sqrt{2/g}\,\sqrt{2g}=2$,
\[
\boxed{\frac{T_{\rm line}}{T_{\rm cyc}}
=\frac2\pi\sqrt{1+\frac{\pi^2}{4}}
=\sqrt{1+\frac4{\pi^2}}}\approx1.185447.
\]


\Needspace{8\baselineskip}
\subsection*{22.2 More General Variations}

\Needspace{5\baselineskip}
\par\bigskip
\noindent\textbf{22.2.1}\quad
\erratum{Take $m>0$ and fixed positions at fixed times $t_1<t_2$.}
\begin{enumerate}
\item[(a)] Develop the equations of motion corresponding to
$L=\tfrac{1}{2}m(\dot{x}^{2}+\dot{y}^{2})$.
\item[(b)] In what sense do your solutions minimize the integral
$\displaystyle\int_{t_{1}}^{t_{2}}L\,dt$?
\end{enumerate}
Compare the result for your solution with $x=$ constant,
$y=$ constant.
\solution
(a) Since $L$ is independent of $x,y$,
\[
\frac d{dt}\frac{\partial L}{\partial\dot x}
-\frac{\partial L}{\partial x}=m\ddot x=0,\qquad
m\ddot y=0.
\]
Thus $x=A+Bt$, $y=C+Dt$. For fixed endpoint positions
$\mathbf r_1,\mathbf r_2$ and times $t_1<t_2$, the unique path is
\[
\boxed{\mathbf r_*(t)=\mathbf r_1+
\frac{t-t_1}{t_2-t_1}(\mathbf r_2-\mathbf r_1)}.
\]

(b) Write another admissible path as
$\mathbf r=\mathbf r_*+\boldsymbol\eta$, with
$\boldsymbol\eta(t_1)=\boldsymbol\eta(t_2)=0$. Its action is
\[
S[\mathbf r]=S[\mathbf r_*]
+m\dot{\mathbf r}_*\cdot\int_{t_1}^{t_2}\dot{\boldsymbol\eta}\,dt
+\frac m2\int_{t_1}^{t_2}|\dot{\boldsymbol\eta}|^2dt.
\]
The middle term is zero because the velocity is constant and the
endpoint variation is zero. Therefore
\[
\boxed{S[\mathbf r]-S[\mathbf r_*]
=\frac m2\int_{t_1}^{t_2}|\dot{\boldsymbol\eta}|^2dt\ge0},
\qquad
S[\mathbf r_*]=\frac{m|\mathbf r_2-\mathbf r_1|^2}{2(t_2-t_1)}.
\]
Equality requires the same path. A constant path has zero action,
but is an admissible competitor only if the prescribed endpoint
positions coincide; in that case it is precisely the minimizing solution.


\Needspace{5\baselineskip}
\par\bigskip
\noindent\textbf{22.2.2}\quad From the Lagrangian equations of
motion, Eq.~(22.38), show that a system in stable equilibrium has
a minimum potential energy.
\erratum{Assume a natural mechanical Lagrangian with positive-definite
kinetic energy and a nondegenerate equilibrium.}
\solution
For a natural mechanical system take
$T=\tfrac12M_{ij}(q)\dot q_i\dot q_j$ with positive-definite
mass matrix, and $L=T-V(q)$. At a rest equilibrium $q=q_0$,
$\dot q=\ddot q=0$. The Lagrange equations reduce to
\[
\boxed{\frac{\partial V}{\partial q_i}(q_0)=0}.
\]
Write $q=q_0+\eta$ and
$K_{ij}=\partial_i\partial_jV(q_0)$. To first order,
\[
M(q_0)\ddot\eta+K\eta=0.
\]
Choose normal coordinates by diagonalizing the real symmetric matrix
$M^{-1/2}KM^{-1/2}$. Each coordinate satisfies
\[
\ddot Q_j+\omega_j^2Q_j=0.
\]
A negative $\omega_j^2=-\kappa^2$ yields
$Q_j=Ae^{\kappa t}+Be^{-\kappa t}$ and hence instability under
arbitrarily small perturbations. In the nondegenerate stable case,
all $\omega_j^2>0$, so $K$ is positive definite. Taylor expansion gives
\[
V(q_0+\eta)=V(q_0)+\frac12\eta^TK\eta+O(|\eta|^3)>V(q_0)
\]
for sufficiently small nonzero $\eta$. Thus the equilibrium has a
strict local minimum of potential energy.

\erratum{The usual conclusion assumes a natural mechanical system
with positive kinetic energy. Linear stability alone gives only a
positive-semidefinite Hessian when zero modes occur; higher-order terms
then decide the result. Lagrange's equations do not turn stationarity
by itself into a minimum.}
For example, $V(q)=q^4$ has a stable minimum despite $V''(0)=0$,
whereas $V(q)=-q^4$ has the same zero Hessian and is unstable.
Conversely, a strict potential well together with positive kinetic
energy confines sufficiently small-energy motions: energy conservation
prevents them from reaching a nearby boundary on which $V$ is higher.
This is the energy reason for stable small oscillations about a minimum.


\Needspace{5\baselineskip}
\par\bigskip
\noindent\textbf{22.2.3}\quad Write out the Lagrangian equations
of motion of a particle in spherical coordinates for potential
$V$ equal to a constant. Identify the terms corresponding to
(a) centrifugal force and (b) Coriolis force.
\erratum{Write the coordinate equations on $r>0$, $0<\theta<\pi$.}
\solution
In spherical coordinates, with $V=V_0$,
\[
L=\frac m2\left(\dot r^2+r^2\dot\theta^2+
r^2\sin^2\theta\,\dot\varphi^2\right)-V_0.
\]
For $r$ the derivatives are
$L_{\dot r}=m\dot r$ and
$L_r=mr(\dot\theta^2+\sin^2\theta\dot\varphi^2)$.
For $\theta$ they are
$L_{\dot\theta}=mr^2\dot\theta$ and
$L_\theta=mr^2\sin\theta\cos\theta\dot\varphi^2$.
For $\varphi$, $L_\varphi=0$ and
$L_{\dot\varphi}=mr^2\sin^2\theta\dot\varphi$.
Differentiating the momenta and dividing by the appropriate scale
factors gives
\[
\boxed{\begin{aligned}
\ddot r&=r\dot\theta^2+r\sin^2\theta\,\dot\varphi^2,\\
r\ddot\theta&=r\sin\theta\cos\theta\,\dot\varphi^2
-2\dot r\dot\theta,\\
r\sin\theta\,\ddot\varphi
&=-2\dot r\sin\theta\,\dot\varphi
-2r\cos\theta\,\dot\theta\dot\varphi.
\end{aligned}}
\]
Equivalently, $mr^2\sin^2\theta\dot\varphi$ is conserved.

The squared-angular-speed terms on the right in the first two
equations are the centrifugal terms of the moving spherical frame.
The products with the factor two are its Coriolis-type terms:
$-2m\dot r\dot\theta$, $-2m\dot r\sin\theta\dot\varphi$,
and $-2mr\cos\theta\dot\theta\dot\varphi$.
There is no physical external force for constant $V$; these terms
account for the changing coordinate basis. Moving them to the left
recovers the three components of the zero Cartesian acceleration.


\Needspace{5\baselineskip}
\par\bigskip
\noindent\textbf{22.2.4}\quad The spherical pendulum consists of a
mass on a wire of length $l$, free to move in polar angle
$\theta$ and azimuth angle $\varphi$ (Fig.~22.8).
\erratum{Measure $\theta$ from the downward vertical.}
\begin{enumerate}
\item[(a)] Set up the Lagrangian for this physical system.
\item[(b)] Develop the Lagrangian equations of motion.
\end{enumerate}
\solution
As in Fig.~22.8, measure $\theta$ from the downward vertical.
The height relative to the pivot is $z=-l\cos\theta$, and the
speed squared is
$l^2(\dot\theta^2+\sin^2\theta\dot\varphi^2)$.
Thus, apart from an irrelevant constant,
\[
\boxed{L=\frac{ml^2}{2}
(\dot\theta^2+\sin^2\theta\dot\varphi^2)+mgl\cos\theta}.
\]
The polar derivatives are
\[
L_{\dot\theta}=ml^2\dot\theta,\qquad
L_\theta=ml^2\sin\theta\cos\theta\dot\varphi^2-mgl\sin\theta.
\]
Their Euler equation is
\[
\boxed{\ddot\theta-\sin\theta\cos\theta\,\dot\varphi^2
+\frac gl\sin\theta=0}.
\]
The azimuth is cyclic, so
\[
\boxed{\frac d{dt}(ml^2\sin^2\theta\,\dot\varphi)=0}.
\]
Where $\sin\theta\ne0$, the latter expands to
\[
\ddot\varphi+2\cot\theta\,\dot\theta\dot\varphi=0.
\]
The conserved momentum is the vertical angular momentum.
At a polar coordinate singularity the conserved-momentum form is the
appropriate one; division by $\sin\theta$ is unnecessary.


\Needspace{5\baselineskip}
\par\bigskip
\noindent\textbf{22.2.5}\quad Show that the Lagrangian
\[
L=m_{0}c^{2}\left(1-\sqrt{1-\frac{v^{2}}{c^{2}}}\right)-V(r)
\]
leads to a relativistic form of Newton's second law of motion,
\[
\frac{d}{dt}\left(\frac{m_{0}v_{i}}{\sqrt{1-v^{2}/c^{2}}}\right)
=F_{i},
\]
in which the force components are $F_{i}=-\partial V/\partial
x_{i}$.
\erratum{Assume $m_0>0$ and $v<c$.}
\solution
Let $v^2=\sum_j\dot x_j^2<c^2$. Since
$\partial v^2/\partial\dot x_i=2\dot x_i$,
\[
\frac{\partial L}{\partial\dot x_i}
=-m_0c^2\frac{1}{2\sqrt{1-v^2/c^2}}
\left(-\frac{2\dot x_i}{c^2}\right)
=\frac{m_0v_i}{\sqrt{1-v^2/c^2}}.
\]
The constant rest-energy term has zero derivative, and the spatial
derivative is $\partial L/\partial x_i=-\partial_iV$.
Consequently
\[
\boxed{\frac d{dt}
\left(\frac{m_0v_i}{\sqrt{1-v^2/c^2}}\right)
=-\partial_iV=F_i}.
\]
Writing $\gamma=(1-v^2/c^2)^{-1/2}$ gives
$\dot\gamma=\gamma^3(\mathbf v\cdot\mathbf a)/c^2$, so the
same equation is
\[
m_0\gamma a_i+
\frac{m_0\gamma^3}{c^2}v_i(\mathbf v\cdot\mathbf a)=F_i.
\]
In the limit $v/c\to0$, this reduces to $m_0a_i=F_i$.


\Needspace{5\baselineskip}
\par\bigskip
\noindent\textbf{22.2.6}\quad The Lagrangian for a particle with
charge $q$ in an electromagnetic field described by scalar
potential $\varphi$ and vector potential $\mathbf{A}$ is
\[
L=\frac{1}{2}mv^{2}-q\varphi+q\mathbf{A}\cdot\mathbf{v}.
\]
Find the equation of motion of the charged particle.
\par\emph{Hint.}~$(d/dt)A_{j}=\partial A_{j}/\partial t
+\displaystyle\sum_{i}(\partial A_{j}/\partial x_{i})\dot{x}_{i}$.
The dependence of the force fields $\mathbf{E}$ and $\mathbf{B}$
on the potentials $\varphi$ and $\mathbf{A}$ is developed in
Section~3.9; see in particular Eq.~(3.108).
\par\nobreak\hfill\emph{ANS.}\enspace
$m\ddot{x}_{i}=q[\mathbf{E}+\mathbf{v}\times\mathbf{B}]_{i}$.
\solution
Treat the potentials as functions of both position and time. The
canonical momentum and spatial derivative are
\[
\frac{\partial L}{\partial\dot x_i}=m\dot x_i+qA_i,\qquad
\frac{\partial L}{\partial x_i}
=-q\,\partial_i\varphi+q\sum_j\dot x_j\partial_iA_j.
\]
Along the trajectory,
\[
\frac{dA_i}{dt}=\partial_tA_i+\sum_j\dot x_j\partial_jA_i.
\]
The Lagrange equation therefore gives
\[
m\ddot x_i
=q(-\partial_i\varphi-\partial_tA_i)
+q\sum_j\dot x_j(\partial_iA_j-\partial_jA_i).
\]
Define $\mathbf E=-\nabla\varphi-\partial_t\mathbf A$ and
$\mathbf B=\nabla\times\mathbf A$.
For example, the last sum at $i=x$ is
\[
v_y(\partial_xA_y-\partial_yA_x)
+v_z(\partial_xA_z-\partial_zA_x)
=v_yB_z-v_zB_y.
\]
The other two components follow by cyclic permutation. Thus
\[
\boxed{m\ddot{\mathbf r}
=q(\mathbf E+\mathbf v\times\mathbf B)}.
\]
The derivatives acting on the position-dependent vector potential are
essential; omitting them would omit the magnetic force.


\Needspace{5\baselineskip}
\par\bigskip
\noindent\textbf{22.2.7}\quad Consider a system in which the
Lagrangian is given by
\[
L(q_{i},\dot{q}_{i})=T(q_{i},\dot{q}_{i})-V(q_{i}),
\]
where $q_{i}$ and $\dot{q}_{i}$ represent sets of variables. The
potential energy $V$ is independent of velocity and neither $T$
nor $V$ has any explicit time dependence.
\begin{enumerate}
\item[(a)] Show that
\[
\frac{d}{dt}\left(\sum_{j}\dot{q}_{j}\,
\frac{\partial L}{\partial\dot{q}_{j}}-L\right)=0.
\]
\item[(b)] The constant quantity
\[
\sum_{j}\dot{q}_{j}\,\frac{\partial L}{\partial\dot{q}_{j}}-L
\]
defines the Hamiltonian $H$. Show that under the preceding
assumed conditions, $H$ satisfies $H=T+V$, and is therefore the
total energy.
\end{enumerate}
\par\emph{Note.}~The kinetic energy $T$ is a \erratum{homogeneous} quadratic function
of the $\dot{q}_{i}$.
\solution
(a) Put $p_j=\partial L/\partial\dot q_j$ and
$H=\sum_j\dot q_jp_j-L$. Its full derivative is
\[
\begin{aligned}
\frac{dH}{dt}
&=\sum_j(\ddot q_jp_j+\dot q_j\dot p_j)
-\sum_j\left(L_{q_j}\dot q_j+L_{\dot q_j}\ddot q_j\right)
-L_t\\
&=\sum_j\dot q_j(\dot p_j-L_{q_j})-L_t.
\end{aligned}
\]
The Lagrange equations make the sum zero, and the stated absence of
explicit time dependence makes $L_t=0$. Hence
\[
\boxed{\frac{dH}{dt}=0}.
\]

(b) Use the note's homogeneous quadratic kinetic energy:
\[
T=\frac12\sum_{i,j}M_{ij}(q)\dot q_i\dot q_j,
\qquad M_{ij}=M_{ji}.
\]
Then
\[
\sum_j\dot q_j\frac{\partial T}{\partial\dot q_j}
=\sum_{i,j}M_{ij}\dot q_i\dot q_j=2T.
\]
Since $V$ has no velocity dependence, $L_{\dot q_j}=T_{\dot q_j}$,
and therefore
\[
\boxed{H=2T-(T-V)=T+V}.
\]
The quadratic homogeneity is used specifically in (b); conservation
of the Legendre-transform expression in (a) does not require it.


\Needspace{5\baselineskip}
\par\bigskip
\noindent\textbf{22.2.8}\quad The Lagrangian for a vibrating
string (small-amplitude vibrations) is
\[
L=\int\left(\frac{1}{2}\rho u_{t}^{2}
-\frac{1}{2}\tau u_{x}^{2}\right)dx,
\]
where $\rho$ is the (constant) linear mass density and $\tau$ is
the (constant) tension. The $x$-integration is over the length
of the string. Show that application of Hamilton's principle to
the Lagrangian density (the integrand), now with two independent
variables, leads to the classical wave equation,
\[
\frac{\partial^{2}u}{\partial x^{2}}
=\frac{\rho}{\tau}\,\frac{\partial^{2}u}{\partial t^{2}}.
\]
\erratum{Assume $\rho,\tau>0$ and variations vanishing on the
spatial and temporal boundaries.}
\solution
Hamilton's principle applies to the action
\[
S[u]=\int_{t_1}^{t_2}dt\int_{x_1}^{x_2}dx\,
\left(\frac12\rho u_t^2-\frac12\tau u_x^2\right).
\]
Let $\delta u=\eta$ vanish at the fixed spatial and temporal
boundaries. Then
\[
\delta S=\iint(\rho u_t\eta_t-\tau u_x\eta_x)\,dx\,dt.
\]
Integrating the first term by parts in $t$ and the second in $x$,
all boundary terms vanish and
\[
\delta S=\iint(-\rho u_{tt}+\tau u_{xx})\eta\,dx\,dt.
\]
Arbitrariness of the interior variation gives
\[
\boxed{u_{xx}=\frac{\rho}{\tau}u_{tt}},
\]
with wave speed $\sqrt{\tau/\rho}$ for positive tension and density.
The constancy of $\rho,\tau$ permits moving them outside the derivatives.


\Needspace{5\baselineskip}
\par\bigskip
\noindent\textbf{22.2.9}\quad Show that the stationary value of
the total energy of the electrostatic field of Example~22.2.3 is
a minimum.
\erratum{The region is charge-free and connected, the boundary
potential is fixed, and the permittivity is a positive constant.}
\par\emph{Hint.}~Investigate the $\alpha^{2}$ terms of $J$.
\solution
Let $\varphi_*$ be the harmonic potential satisfying the fixed
boundary values in the charge-free region $\Omega$:
$\nabla^2\varphi_*=0$. For any admissible variation
$\varphi=\varphi_*+\alpha\eta$ with $\eta=0$ on the boundary,
the physical field energy is
\[
U[\varphi]=\frac\varepsilon2\int_\Omega|\nabla\varphi|^2dV.
\]
Expansion is exact because the functional is quadratic:
\[
U[\varphi_*+\alpha\eta]-U[\varphi_*]
=\alpha\varepsilon\int_\Omega\nabla\varphi_*\cdot\nabla\eta\,dV
+\frac{\alpha^2\varepsilon}{2}\int_\Omega|\nabla\eta|^2dV.
\]
The cross integral is
\[
\int_{\partial\Omega}\eta\,\partial_n\varphi_*\,dS
-\int_\Omega\eta\,\nabla^2\varphi_*\,dV=0.
\]
Thus
\[
\boxed{U[\varphi_*+\alpha\eta]-U[\varphi_*]
=\frac{\alpha^2\varepsilon}{2}\int_\Omega|\nabla\eta|^2dV\ge0}.
\]
For $\varepsilon>0$ in a connected region with fixed boundary
potential, equality requires $\nabla\eta=0$ and hence $\eta=0$.
This proves a strict global minimum. The book's $J$ omits only the
positive constant $\varepsilon/2$, so the conclusion is identical.


\Needspace{5\baselineskip}
\par\bigskip
\noindent\textbf{22.2.10}\quad The Lagrangian (per unit volume) of
an electromagnetic field with a charge density $\rho$ and current
density $\mathbf{J}$ is given by
\[
L=\frac{1}{2}\left(\varepsilon_{0}E^{2}
-\frac{1}{\mu_{0}}B^{2}\right)-\rho\varphi
+\mathbf{J}\cdot\mathbf{A}.
\]
Show that Lagrange's equations lead to two of Maxwell's
equations. (The remaining two are a consequence of the definition
of $\mathbf{E}$ and $\mathbf{B}$ in terms of $\mathbf{A}$ and
$\varphi$.)
\par\emph{Hint.}~Take $\varphi$ and the components of
$\mathbf{A}$ as dependent variables; and $x$, $y$, $z$, and $t$
as independent variables. $\mathbf{E}$ and $\mathbf{B}$ are given
in terms of $\mathbf{A}$ and $\varphi$ by Eq.~(3.108).
\erratum{Use prescribed sources satisfying
$\partial_t\rho+\nabla\cdot\mathbf J=0$ and variations vanishing
on the spacetime boundary.}
\solution
Use
\[
\mathbf E=-\nabla\varphi-\partial_t\mathbf A,\qquad
\mathbf B=\nabla\times\mathbf A,
\]
with prescribed sources $\rho,\mathbf J$. The action is the spacetime
integral of the given density. Its variation is
\[
\delta S=\int dt\,d^3x\left[
\varepsilon_0\mathbf E\cdot\delta\mathbf E
-\frac1{\mu_0}\mathbf B\cdot\delta\mathbf B
-\rho\,\delta\varphi+\mathbf J\cdot\delta\mathbf A\right].
\]

First vary only $\varphi$. Then $\delta\mathbf E=-\nabla\delta\varphi$
and $\delta\mathbf B=0$. Spatial integration by parts gives
\[
\delta_\varphi S=\int
(\varepsilon_0\nabla\cdot\mathbf E-\rho)\delta\varphi\,dt\,d^3x.
\]
Boundary variations vanish. Hence
\[
\boxed{\nabla\cdot\mathbf E=\rho/\varepsilon_0}.
\]
In component Euler notation, this uses
$L_\varphi=-\rho$ and $L_{\partial_i\varphi}=-\varepsilon_0E_i$.

Next vary only $\mathbf A$. Now
$\delta\mathbf E=-\partial_t\delta\mathbf A$ and
$\delta\mathbf B=\nabla\times\delta\mathbf A$. The time term becomes
\[
-\int\varepsilon_0\mathbf E\cdot\partial_t\delta\mathbf A
=\int\varepsilon_0\partial_t\mathbf E\cdot\delta\mathbf A.
\]
For the magnetic term, the identity
\[
\nabla\cdot(\delta\mathbf A\times\mathbf B)
=\mathbf B\cdot(\nabla\times\delta\mathbf A)
-\delta\mathbf A\cdot(\nabla\times\mathbf B)
\]
shows that its volume integral becomes
$-\mu_0^{-1}\int(\nabla\times\mathbf B)\cdot\delta\mathbf A$.
The three independent components therefore give
\[
\mathbf J+\varepsilon_0\partial_t\mathbf E
-\mu_0^{-1}\nabla\times\mathbf B=0,
\]
or
\[
\boxed{\nabla\times\mathbf B
=\mu_0\mathbf J+\mu_0\varepsilon_0\partial_t\mathbf E}.
\]
Finally the definitions themselves give
$\nabla\cdot\mathbf B=\nabla\cdot(\nabla\times\mathbf A)=0$ and
$\nabla\times\mathbf E=-\partial_t(\nabla\times\mathbf A)
=-\partial_t\mathbf B$, the remaining two Maxwell equations.


\Needspace{8\baselineskip}
\subsection*{22.3 Constrained Minima/Maxima}

\Needspace{5\baselineskip}
\par\bigskip
\noindent\textbf{22.3.0}\quad The following problems are to be
solved by using Lagrangian multipliers.
\solution
This numbered entry is the source's instruction for the following
problems. For an objective $F(q)$ constrained by $G(q)=C$,
a permissible displacement satisfies $\nabla G\cdot\delta q=0$.
At a constrained stationary point, $\nabla F\cdot\delta q=0$
for every such tangent displacement. If $\nabla G\ne0$, the two
normal vectors must therefore be proportional:
\[
\boxed{\nabla F-\lambda\nabla G=0,\qquad G=C}.
\]
For several independent constraints, $\nabla F$ is a linear combination
of their gradients. The multiplier equations find stationary candidates;
one must still examine boundary behavior or use an inequality to decide
whether each candidate is a maximum, minimum, or neither. The following
solutions include that classification.


\Needspace{5\baselineskip}
\par\bigskip
\noindent\textbf{22.3.1}\quad The ground-state energy of a quantum
particle of mass $m$ in a pillbox (right-circular cylinder) is
given by
\[
E=\frac{\hbar^{2}}{2m}
\left(\frac{(2.4048)^{2}}{R^{2}}+\frac{\pi^{2}}{H^{2}}\right),
\]
in which $R$ is the radius and $H$ is the height of the pillbox.
Find the ratio of $R$ to $H$ that will minimize the energy for a
fixed positive volume.
\erratum{Here $2.4048$ is the stated rounded approximation to the
first positive zero of $J_0$.}
\solution
Put $j=2.4048$, retain the positive constant $\hbar^2/(2m)$
outside the minimization, and impose $\pi R^2H=V$.
For
\[
\mathcal F=\frac{j^2}{R^2}+\frac{\pi^2}{H^2}
+\lambda(\pi R^2H-V),
\]
the two multiplier equations are
\[
-\frac{2j^2}{R^3}+2\lambda\pi RH=0,\qquad
-\frac{2\pi^2}{H^3}+\lambda\pi R^2=0.
\]
Hence
\[
\lambda=\frac{j^2}{\pi R^4H}
=\frac{2\pi}{R^2H^3},
\qquad j^2H^2=2\pi^2R^2.
\]
The requested ratio is
\[
\boxed{\frac RH=\frac{2.4048}{\pi\sqrt2}
=0.541270169274\ldots}.
\]
To classify it, eliminate $H=V/(\pi R^2)$:
\[
\frac{2mE}{\hbar^2}
=\frac{j^2}{R^2}+\frac{\pi^4R^4}{V^2}.
\]
This diverges as $R\downarrow0$ or $R\to\infty$, and its derivative
vanishes only at $R^6=j^2V^2/(2\pi^4)$.
Thus the stationary ratio gives the unique global minimum.
\erratum{The manual replaces the supplied $2.4048$ by $2.048$
in both its derivative and final ratio. The coefficient in the
original problem is $2.4048$, which is used throughout above.}


\Needspace{5\baselineskip}
\par\bigskip
\noindent\textbf{22.3.2}\quad \erratum{Under the textbook's stipulated postal rule,}
the total length plus girth may not exceed 36 inches. Using Lagrange multipliers, find the dimensions of the
rectangular parallelepiped of maximum volume subject to this
constraint.
\solution
Use the stated 36-inch mathematical constraint. Let the length be
$L$ and the two transverse widths be $w,h$. A maximum uses the
whole allowance: otherwise increasing $L$ raises the volume while
remaining admissible. Thus
\[
V=Lwh,\qquad L+2w+2h=36.
\]
For $\mathcal F=Lwh-\lambda(L+2w+2h-36)$,
\[
wh=\lambda,\qquad Lh=2\lambda,\qquad Lw=2\lambda.
\]
All dimensions are positive, so division by $wh$ gives
$L=2w=2h$. The constraint becomes $L+L+L=36$, hence
\[
\boxed{L=12\ {\rm in},\qquad w=h=6\ {\rm in},\qquad
V_{\max}=432\ {\rm in}^3}.
\]
For a global check, the three positive numbers $L,2w,2h$ have
fixed sum 36, and their product is $4V$. The arithmetic--geometric
mean inequality gives
\[
4V\le(36/3)^3=1728,
\]
with equality exactly at $L=2w=2h$. The identified length is also
the longest side, consistently with the length-plus-girth convention.


\Needspace{5\baselineskip}
\par\bigskip
\noindent\textbf{22.3.3}\quad A thermal nuclear reactor is subject
to the constraint
\[
\varphi(a,b,c)=\left(\frac{\pi}{a}\right)^{2}
+\left(\frac{\pi}{b}\right)^{2}+\left(\frac{\pi}{c}\right)^{2}
=B^{2},\quad\text{a constant},
\]
where the reactor is a rectangular parallelepiped of sides $a$,
$b$, and $c$. \erratum{Take $a,b,c,B>0$.} Find the ratios of $a$, $b$, and $c$ that \erratum{minimize}
the reactor volume. \erratum{Show also that no maximum exists.}
\par\nobreak\hfill\emph{ANS.}\enspace$a=b=c$, cube.
\solution
For positive side lengths, extremizing $abc$ is equivalent to
extremizing $\log(abc)$. Apply a multiplier to
\[
\mathcal F=\log a+\log b+\log c+
\lambda\left[\pi^2(a^{-2}+b^{-2}+c^{-2})-B^2\right].
\]
The three equations are
\[
\frac1a-\frac{2\lambda\pi^2}{a^3}=0,\qquad
\frac1b-\frac{2\lambda\pi^2}{b^3}=0,\qquad
\frac1c-\frac{2\lambda\pi^2}{c^3}=0.
\]
Thus $a^2=b^2=c^2=2\lambda\pi^2$ and positivity yields
$a=b=c$. The constraint gives
\[
a=b=c=\frac{\sqrt3\pi}{B}.
\]

\erratum{The cube is the \emph{minimum}-volume reactor under this
constraint. The printed request for a maximum is false: the volume
has no finite upper bound.}
To prove both statements, set $u=a^{-2}$, $v=b^{-2}$,
$w=c^{-2}$, so $u+v+w=B^2/\pi^2$. Then
\[
uvw\le\left(\frac{B^2}{3\pi^2}\right)^3,\qquad
abc=(uvw)^{-1/2}\ge\boxed{\frac{3\sqrt3\pi^3}{B^3}}.
\]
Equality occurs only at $u=v=w$, i.e. the cube.

For an explicit unbounded family, take
\[
a=b=L,\qquad
c=\frac{\pi}{\sqrt{B^2-2\pi^2/L^2}},
\qquad L>\frac{\sqrt2\pi}{B}.
\]
It satisfies the original constraint exactly, while
$abc=\pi L^2/\sqrt{B^2-2\pi^2/L^2}\to\infty$.
This resolves the original maximum question as well as identifying
the correctly classified stationary cube.


\Needspace{5\baselineskip}
\par\bigskip
\noindent\textbf{22.3.4}\quad For a lens of focal length $f$, the
object distance $p$ and the image distance $q$ are related by
$1/p+1/q=1/f$. Find the minimum object-image distance $(p+q)$
for fixed $f$. Assume real object and image ($p$ and $q$ both
positive).
\solution
The assumptions $p,q>0$ imply $f>0$. Use
\[
\mathcal F=p+q+\lambda\left(\frac1p+\frac1q-\frac1f\right).
\]
Differentiation gives $1-\lambda/p^2=0$ and
$1-\lambda/q^2=0$. Thus $p^2=q^2$, so $p=q$.
The lens equation then gives $p=q=2f$ and
\[
\boxed{(p+q)_{\min}=4f}.
\]
For a global proof, put $S=p+q$. The constraint implies $pq=fS$.
Since $(p-q)^2\ge0$,
\[
S^2\ge4pq=4fS.
\]
Divide by the positive $S$ to obtain $S\ge4f$, with equality only
for the multiplier solution.


\Needspace{5\baselineskip}
\par\bigskip
\noindent\textbf{22.3.5}\quad You have an ellipse
$(x/a)^{2}+(y/b)^{2}=1$. Find the inscribed rectangle of maximum
area. Show that the ratio of the area of the maximum-area
rectangle to the area of the ellipse is $2/\pi=0.6366$.
\solution
First take the rectangle's vertices to be $(\pm x,\pm y)$ with
$x,y>0$. Its area is $A=4xy$ and its constraint is
$x^2/a^2+y^2/b^2=1$. For
\[
\mathcal F=4xy-\lambda(x^2/a^2+y^2/b^2-1),
\]
the equations are
\[
4y=\frac{2\lambda x}{a^2},\qquad
4x=\frac{2\lambda y}{b^2}.
\]
Multiplying respectively by $x,y$ gives
$x^2/a^2=y^2/b^2$. Each is therefore $1/2$, so
\[
\boxed{x=\frac a{\sqrt2},\quad y=\frac b{\sqrt2},\quad
A_{\max}=2ab}.
\]
The side lengths are $\sqrt2a,\sqrt2b$ and
\[
\boxed{\frac{A_{\max}}{A_{\rm ellipse}}
=\frac{2ab}{\pi ab}=\frac2\pi=0.6366197724\ldots}.
\]

This is also the maximum if the rectangle's orientation is unrestricted.
Let its center be $h$, its orthonormal directions be $e_1,e_2$,
and its half-lengths be $u,v$. For the ellipse matrix
$Q=\operatorname{diag}(a^{-2},b^{-2})$, averaging the four vertex
inequalities gives
\[
h^TQh+u^2q_1+v^2q_2\le1,\qquad q_i=e_i^TQe_i.
\]
Thus $2uv\sqrt{q_1q_2}\le1$. In this rotated basis,
$q_1q_2=\det Q+(e_1^TQe_2)^2\ge1/(a^2b^2)$,
and hence $4uv\le2ab$. The aligned multiplier solution attains this
bound, proving the global claim.


\Needspace{5\baselineskip}
\par\bigskip
\noindent\textbf{22.3.6}\quad A rectangular parallelepiped is
inscribed in an ellipsoid of semiaxes $a$, $b$, and $c$. Maximize
the volume of the inscribed rectangular parallelepiped. Show
that the ratio of the maximum volume to the volume of the
ellipsoid is $\erratum{\dfrac{2}{\pi\sqrt{3}}}\approx0.367$.
\solution
For an aligned box let the positive half-lengths be $x,y,z$.
Maximize $V=8xyz$ subject to
$x^2/a^2+y^2/b^2+z^2/c^2=1$. The multiplier equations give
\[
8yz=\frac{2\lambda x}{a^2},\quad
8xz=\frac{2\lambda y}{b^2},\quad
8xy=\frac{2\lambda z}{c^2}.
\]
Multiplying by $x,y,z$ respectively makes all left sides equal.
Thus $x^2/a^2=y^2/b^2=z^2/c^2=1/3$, giving
\[
\boxed{x=\frac a{\sqrt3},\quad y=\frac b{\sqrt3},\quad
z=\frac c{\sqrt3},\qquad V_{\max}=\frac{8abc}{3\sqrt3}}.
\]
Dividing by the ellipsoid's $4\pi abc/3$ yields
\[
\boxed{\frac{V_{\max}}{V_{\rm ellipsoid}}
=\frac{2}{\pi\sqrt3}=0.3675525969\ldots}.
\]
\erratum{The inline ratio must be read as $2/(\pi\sqrt3)$;
$(2/\pi)\sqrt3$ would exceed one and would not match the stated
decimal approximation.}

For completeness, the same bound covers arbitrary orientations and
centers. With $Q=\operatorname{diag}(a^{-2},b^{-2},c^{-2})$,
orthonormal box directions $e_i$ and half-lengths $u_i$, the average
of the eight vertex inequalities gives
$\sum_iq_i u_i^2\le1$, where $q_i=e_i^TQe_i$; the nonnegative
center term can be discarded. The arithmetic--geometric mean bound
then gives $\prod_i(q_i u_i^2)\le1/27$.
The Gram determinant inequality gives
$\prod_iq_i\ge\det Q=1/(a^2b^2c^2)$: the squared volume of
the vectors $Q^{1/2}e_i$ is at most the product of their squared
lengths. Therefore
\[
8u_1u_2u_3\le\frac8{\sqrt{27}\sqrt{q_1q_2q_3}}
\le\frac{8abc}{3\sqrt3}.
\]
The aligned box attains equality.


\Needspace{5\baselineskip}
\par\bigskip
\noindent\textbf{22.3.7}\quad Find the maximum value of the
directional derivative of $\varphi(x,y,z)$,
\erratum{a real differentiable function at a fixed point,}
\[
\frac{d\varphi}{ds}=\frac{\partial\varphi}{\partial x}\cos\alpha
+\frac{\partial\varphi}{\partial y}\cos\beta
+\frac{\partial\varphi}{\partial z}\cos\gamma,
\]
subject to the constraint,
\[
\cos^{2}\alpha+\cos^{2}\beta+\cos^{2}\gamma=1.
\]
\solution
At the point in question write $\mathbf g=\nabla\varphi$ and
$\mathbf n=(\cos\alpha,\cos\beta,\cos\gamma)$.
Maximize $D=\mathbf g\cdot\mathbf n$ subject to
$\mathbf n\cdot\mathbf n=1$. The multiplier equations for
$D-\lambda(\mathbf n\cdot\mathbf n-1)$ are
\[
g_i-2\lambda n_i=0.
\]
If $\mathbf g\ne0$, normalization gives
$4\lambda^2=|\mathbf g|^2$, and hence
\[
\mathbf n_\pm=\pm\frac{\mathbf g}{|\mathbf g|},\qquad
D_\pm=\pm|\mathbf g|.
\]
The maximum is therefore
\[
\boxed{\left(\frac{d\varphi}{ds}\right)_{\max}
=|\nabla\varphi|,\qquad
\mathbf n_{\max}=\frac{\nabla\varphi}{|\nabla\varphi|}}.
\]
Cauchy--Schwarz proves the global bound
$\mathbf g\cdot\mathbf n\le|\mathbf g||\mathbf n|=|\mathbf g|$.
If $\nabla\varphi=0$, every unit direction gives zero, so every
direction is a maximizer and no preferred direction is defined.


\Needspace{8\baselineskip}
\subsection*{22.4 Variation with Constraints}

\Needspace{5\baselineskip}
\par\bigskip
\noindent\textbf{22.4.1}\quad A particle of mass $m$ is on a
frictionless horizontal surface. In terms of plane polar
coordinates $(r,\theta)$, it is constrained to move so that
$\theta=\omega t$ (accomplished by pushing it with a rotating
radial arm against which it can slide frictionlessly). With the
initial conditions
\[
t=0,\qquad r=r_{0},\qquad\dot{r}=0:
\]
\begin{enumerate}
\item[(a)] Find the radial position as a function of time.
\par\nobreak\hfill\emph{ANS.}\enspace$r(t)=r_{0}\cosh\omega t$.
\item[(b)] Find the force exerted on the particle by the
constraint.
\par\nobreak\hfill\emph{ANS.}\enspace
$F^{(c)}=2m\dot{r}\omega=2mr_{0}\omega^{2}\sinh\omega t$.
\end{enumerate}
\solution
Use polar coordinates with the positive tangential direction
$\mathbf e_\theta$.  Before imposing the constraint,
\[
 L=\frac m2(\dot r^2+r^2\dot\theta^2),\qquad
 \phi=\theta-\omega t=0.
\]
With the convention $L+\lambda\phi$, the two Lagrange equations are
\[
 \frac{d}{dt}(m\dot r)-mr\dot\theta^2=0,\qquad
 \frac{d}{dt}(mr^2\dot\theta)=\lambda.
\]
The multiplier in the second equation is a torque, since its
coordinate is an angle.

(a) Substitute $\dot\theta=\omega$ into the first equation:
\[
 \ddot r-\omega^2r=0,\qquad
 r=A\cosh(\omega t)+B\sinh(\omega t).
\]
The initial conditions give $A=r_0$, $B=0$, so
\[
 \boxed{r(t)=r_0\cosh(\omega t)},\qquad
 \dot r=r_0\omega\sinh(\omega t).
\]
The $\omega=0$ limit is $r=r_0$.

(b) The frictionless arm pushes perpendicular to itself, along
$\mathbf e_\theta$.  Since $rF_\theta=\lambda$,
\[
 \lambda=\omega\frac{d}{dt}(mr^2)=2m\omega r\dot r,\qquad
 \boxed{\mathbf F^{(c)}=2m\omega\dot r\,\mathbf e_\theta
 =2mr_0\omega^2\sinh(\omega t)\,\mathbf e_\theta}.
\]
There is no radial applied force: the radial acceleration component
$\ddot r-r\omega^2$ vanishes.  As an energy check,
\[
 \frac{d}{dt}\frac m2(\dot r^2+r^2\omega^2)
 =m\dot r(\ddot r+r\omega^2)
 =2m\omega^2r\dot r
 =\mathbf F^{(c)}\cdot\mathbf v.
\]
Thus the motor turning the arm supplies the increasing kinetic energy.


\Needspace{5\baselineskip}
\par\bigskip
\noindent\textbf{22.4.2}\quad A point mass $m$ is moving on a
flat, horizontal, frictionless plane. The mass is constrained by
a string to move radially inward at a constant rate. Using plane
polar coordinates $(r,\theta)$, $r=r_{0}-kt$,
\erratum{with $m,r_0,k>0$ and $0\le t<r_0/k$:}
\begin{enumerate}
\item[(a)] Set up the Lagrangian.
\item[(b)] Obtain the constrained Lagrange equations.
\item[(c)] Solve the $\theta$-dependent Lagrange equation to
obtain $\omega(t)$, the angular velocity. What is the physical
significance of the constant of integration that you get from
your ``free'' integration?
\item[(d)] Using the $\omega(t)$ from part \erratum{(c)}, solve the
$r$-dependent (constrained) Lagrange equation to obtain
$\lambda(t)$. In other words, explain what is happening to the
force of constraint as $r\to0$.
\end{enumerate}
\solution
(a) Take $r$ and $\theta$ as independent coordinates initially:
\[
 L=\frac m2(\dot r^2+r^2\dot\theta^2),\qquad
 \phi(r,t)=r-r_0+kt=0.
\]
Here $k>0$, and the prescribed motion applies for $0\le t<r_0/k$.

(b) Using $L+\lambda\phi$ gives
\[
 \boxed{m\ddot r-mr\dot\theta^2=\lambda},\qquad
 \boxed{\frac{d}{dt}(mr^2\dot\theta)=0},\qquad r=r_0-kt.
\]
Because $\partial\phi/\partial r=1$, $\lambda$ is the signed radial
force, positive outward.  The physical string tension is $-\lambda$.

(c) Integrating the cyclic-coordinate equation once,
\[
 mr^2\dot\theta=\ell,\qquad
 \boxed{\omega(t)=\dot\theta=\frac{\ell}{m(r_0-kt)^2}
 =\omega_0\frac{r_0^2}{(r_0-kt)^2}}.
\]
The constant $\ell=mr_0^2\omega_0$ is the conserved angular momentum
about the hole through which the string is drawn.
\erratum{The manual prints the third power in the denominator of
$\omega(t)$. It must be the second power: the printed expression
would make $mr^2\omega=\ell/r$ vary with time.}
The constraint
force is radial and therefore has zero torque about that point.
A second integration, if desired, gives
\[
 \theta(t)-\theta_0
 =\frac{\ell}{m}\int_0^t\frac{d\tau}{(r_0-k\tau)^2}
 =\frac{\ell}{mk}\left(\frac1{r_0-kt}-\frac1{r_0}\right).
\]

(d) Since $\ddot r=0$, the radial equation becomes
\[
 \boxed{\lambda(t)=-mr\omega^2
       =-\frac{\ell^2}{m(r_0-kt)^3}},\qquad
 T_{\rm string}=\frac{\ell^2}{mr^3}.
\]
For $\ell\ne0$, the tension tends to $+\infty$ as $r\to0^+$.
Maintaining the specified radial speed therefore requires an
unbounded pulling force in that limit.
\erratum{The divergence requires nonzero angular momentum.
If $\ell=0$, then $\omega=0$ and $\lambda=0$ throughout the prescribed
inward motion; the endpoint $r=0$ itself is a polar-coordinate
singularity.}


\Needspace{5\baselineskip}
\par\bigskip
\noindent\textbf{22.4.3}\quad A flexible cable is suspended from
two fixed points. The length of the cable is fixed. Find the
curve that will minimize the total gravitational potential
energy of the cable.
\erratum{Assume a uniform cable in uniform gravity, supports at
different horizontal positions, and a cable length greater than
the straight-line distance between the supports.}
\par\nobreak\hfill\emph{ANS.}\enspace Hyperbolic cosine.
\solution
\erratum{The hyperbolic-cosine answer assumes a uniform cable in a
uniform gravitational field.  A nonuniform mass per unit length
generally produces a different curve.}
Let the constant mass per unit length be $\rho$, let $y$ increase
upward, and suppose the endpoints have $x_1<x_2$.  For a graph $y(x)$,
\[
 U=\rho g\int_{x_1}^{x_2}y\sqrt{1+y'^2}\,dx,\qquad
 L_0=\int_{x_1}^{x_2}\sqrt{1+y'^2}\,dx.
\]
Introduce a constant multiplier $\lambda$ and vary
\[
 \int_{x_1}^{x_2}(\rho gy+\lambda)\sqrt{1+y'^2}\,dx.
\]
There is no explicit $x$ dependence.  Its Beltrami first integral is
\[
 (\rho gy+\lambda)\sqrt{1+y'^2}
 -y'\frac{(\rho gy+\lambda)y'}{\sqrt{1+y'^2}}
 =\frac{\rho gy+\lambda}{\sqrt{1+y'^2}}
 =\rho gc.
\]
For the hanging solution $c>0$.  Put $Y=y+\lambda/(\rho g)$:
\[
 Y=c\sqrt{1+Y'^2},\qquad
 \frac{dY}{\sqrt{Y^2-c^2}}=\pm\frac{dx}{c}.
\]
Since $\int dY/\sqrt{Y^2-c^2}=\operatorname{arcosh}(Y/c)$,
joining the two sides of the lowest point gives
\[
 \boxed{y(x)=c\cosh\frac{x-b}{c}+d},\qquad
 d=-\frac{\lambda}{\rho g}.
\]
The three constants are determined by the two endpoint values and
\[
 L_0=c\left[\sinh\frac{x_2-b}{c}
                 -\sinh\frac{x_1-b}{c}\right].
\]
\erratum{The manual repeats $x_2$ in both terms of this length
formula. The second evaluation is at $x_1$; otherwise the displayed
length would be zero.}

To verify that the hanging solution minimizes energy, rather than
merely makes it stationary, parameterize it by arc length
$s\in[0,L_0]$.  There is a constant $s_0$ such that
\[
 s-s_0=c\sinh\frac{x-b}{c},\qquad
 \mathbf T(s)=\rho g\bigl(c,\ s-s_0\bigr),\qquad
 \mathbf r_*'(s)=\frac{\mathbf T(s)}{|\mathbf T(s)|}.
\]
Thus $\mathbf T'=\rho g\,\mathbf e_y$.  Any other cable of the same
length with the same endpoints has an arc-length parameterization
$\mathbf r(s)$ with $|\mathbf r'|=1$.  Integration by parts gives
\begin{align*}
 U[\mathbf r]-U[\mathbf r_*]
 &=\int_0^{L_0}\mathbf T'\cdot(\mathbf r-\mathbf r_*)\,ds\\
 &=-\int_0^{L_0}\mathbf T\cdot(\mathbf r'-\mathbf r_*')\,ds\\
 &=\int_0^{L_0}\bigl(|\mathbf T|-\mathbf T\cdot\mathbf r'\bigr)\,ds
 \ge0,
\end{align*}
where the endpoint term vanishes and the last inequality is
Cauchy--Schwarz.  Equality requires the same tangent everywhere,
hence the same curve with the fixed endpoints.  A curved cable needs
$L_0$ greater than the straight-line distance between the endpoints;
equality gives the taut straight segment, and a shorter cable cannot
join the endpoints.
\erratum{For vertically aligned supports the smooth graph assumption
fails. In the ideal zero-thickness model a slack cable has a limiting
vertical fold.} If its endpoint heights are
$y_1,y_2$ and $L_0>|y_2-y_1|$, its lowest point is
\[
 y_b=\frac{y_1+y_2-L_0}{2},
\]
since $(y_1-y_b)+(y_2-y_b)=L_0$. The same calibration proof applies
with horizontal tension zero and
$\mathbf T=\rho g(s-s_0)\mathbf e_y$: the cable goes downward for
$s<s_0=y_1-y_b$ and upward for $s>s_0$.


\Needspace{5\baselineskip}
\par\bigskip
\noindent\textbf{22.4.4}\quad A fixed volume of water is rotating
in a cylinder with constant angular velocity $\omega$. Find the
curve of the water surface that will minimize the total
potential energy of the water in the combined
gravitational-centrifugal force field.
\erratum{Assume constant density, uniform gravity, a flat-bottomed
cylinder, no overflow, and negligible surface tension. Allow a dry
central region if required by nonnegative water depth.}
\par\nobreak\hfill\emph{ANS.}\enspace Parabola.
\solution
\erratum{In the manual's energy integral the centrifugal term is
missing the square on $\omega$. Also, its equation
$2gry-\omega^2r^3-2\lambda r=0$ gives
$y=(\omega^2r^2+2\lambda)/(2g)$, with a plus sign for its stated
multiplier convention. The derivation below uses $U+\lambda V$.}
Work in the frame rotating with the water.  The potential per unit
mass for gravity plus the centrifugal force is
\[
 \Phi(r,z)=gz-\frac12\omega^2r^2,
\]
since $-\nabla\Phi=-g\mathbf e_z+\omega^2r\mathbf e_r$.
Let the cylinder have radius $R$, a flat bottom $z=0$, and water
height $h(r)$.  The density $\rho$ is constant.  Integrating through
the depth gives
\[
 U[h]=2\pi\rho\int_0^R
       \left(\frac g2h^2-\frac{\omega^2r^2}{2}h\right)r\,dr,
 \qquad
 V=2\pi\int_0^R h(r)r\,dr.
\]
For $h>0$, stationarity of $U+\lambda V$ gives the pointwise equation
\[
 \rho\left(gh-\frac{\omega^2r^2}{2}\right)+\lambda=0.
\]
Consequently the free surface is
\[
 \boxed{z=h(r)=h_0+\frac{\omega^2r^2}{2g}},
\]
a paraboloid of revolution, whose section through the axis is a
parabola.  When the bottom is completely wet, the volume fixes
\[
 V=\pi R^2h_0+\frac{\pi\omega^2R^4}{4g},\qquad
 \boxed{h_0=\frac{V}{\pi R^2}-\frac{\omega^2R^2}{4g}}.
\]

For a volume-preserving perturbation $\eta$, the first-order energy
change vanishes because $\rho(gh-\omega^2r^2/2)$ is constant, and
\[
 U[h+\eta]-U[h]=\frac{\rho g}{2}\int_{\text{bottom}}\eta^2\,dA\ge0.
\]
This proves the minimum.

\erratum{The full-disk formula requires $h_0\ge0$ and no overflow.
If it predicts a negative depth, a dry central region appears.
The physical minimizer is
$h(r)=\max\{0,C+\omega^2r^2/(2g)\}$, with $C$ chosen from the volume;
the wetted free surface remains parabolic.}
For example, if its inner wet radius is $r_d$, then
$h=\omega^2(r^2-r_d^2)/(2g)$ on $r_d\le r\le R$, and direct integration
gives $V=\pi\omega^2(R^2-r_d^2)^2/(4g)$.
The same minimum follows from convexity of the quadratic integrand,
with the inequality $h\ge0$ imposed.


\Needspace{5\baselineskip}
\par\bigskip
\noindent\textbf{22.4.5}\quad
\erratum{Consider simple closed regular plane curves enclosing
positive area.}
\begin{enumerate}
\item[(a)] Show that for a fixed-length perimeter the plane
figure with maximum area is a circle.
\item[(b)] Show that for a fixed planar area the boundary with
minimum perimeter is a circle.
\end{enumerate}
\par\emph{Hint.}~The \erratum{geometric} radius of curvature $R$ is given by
\[
R=\frac{(r^{2}+r_{\theta}^{2})^{3/2}}
{\erratum{\left|rr_{\theta\theta}-2r_{\theta}^{2}-r^{2}\right|}}.
\]
\par\emph{Note.}~The problems of this section, variation subject
to constraints, are often called \emph{isoperimetric}. The term
arose from problems of maximizing area subject to a fixed
perimeter, as in part (a) of this problem.
\solution
Let a regular simple closed curve be positively oriented and write
$\mathbf r(u)=(x(u),y(u))$, $v=\sqrt{x'^2+y'^2}$.  Its area and length
are
\[
 A=\frac12\oint(xy'-yx')\,du,\qquad L=\oint v\,du.
\]
Both parts lead to the same stationary curves: vary $A+\lambda L$
in part (a), or $L+\mu A$ in part (b), with a constant multiplier.
For $A+\lambda L$, the $x$ Euler equation is
\[
 \frac{y'}2-\frac d{du}\left(-\frac y2+\lambda\frac{x'}v\right)=0,
 \quad\hbox{so}\quad
 \lambda\frac d{du}\frac{x'}v=y'.
\]
The $y$ equation similarly gives
\[
 -\frac{x'}2-\frac d{du}\left(\frac x2+\lambda\frac{y'}v\right)=0,
 \quad\hbox{so}\quad
 \lambda\frac d{du}\frac{y'}v=-x'.
\]
Use arc length $s$, so $x_s^2+y_s^2=1$.  Integrating these equations,
\[
 x+\lambda y_s=C_1,\qquad y-\lambda x_s=C_2.
\]
Squaring and adding yields
\[
 (x-C_1)^2+(y-C_2)^2=\lambda^2.
\]
Thus every regular extremal is a circle.

It remains to prove the requested global maximum and minimum.
Translate the curve so $\int_0^L\mathbf r(s)\,ds=0$.  For a
mean-zero $L$-periodic function $f$, its Fourier expansion gives
\[
 \int_0^L f'^2\,ds
 =\frac L2\sum_{n\ge1}\left(\frac{2\pi n}{L}\right)^2(a_n^2+b_n^2)
 \ge\left(\frac{2\pi}{L}\right)^2
       \frac L2\sum_{n\ge1}(a_n^2+b_n^2)
 =\left(\frac{2\pi}{L}\right)^2\int_0^L f^2\,ds.
\]
Apply this to both coordinates.  Because $|\mathbf r_s|=1$,
\[
 \int_0^L|\mathbf r|^2\,ds
 \le\frac{L^2}{4\pi^2}\int_0^L|\mathbf r_s|^2\,ds
 =\frac{L^3}{4\pi^2}.
\]
Cauchy--Schwarz now gives the isoperimetric inequality:
\[
 2A=\int_0^L(xy_s-yx_s)\,ds
 \le\left(\int_0^L|\mathbf r|^2\,ds\right)^{1/2}
      \left(\int_0^L|\mathbf r_s|^2\,ds\right)^{1/2}
 \le\frac{L^2}{2\pi}.
\]
Equality holds for a circle, for which $\mathbf r_s$ is a constant
multiple of the $90^\circ$ rotation of $\mathbf r$.  Therefore
\[
 \boxed{A_{\max}=\frac{L^2}{4\pi}\quad\text{at fixed }L},
 \qquad
 \boxed{L_{\min}=2\sqrt{\pi A}\quad\text{at fixed }A}.
\]
\erratum{The original textbook hint's expression is a signed radius of curvature with
the opposite sign to the usual counterclockwise convention: for
$r=\text{constant}>0$ it gives $R=-r$.  A geometric radius requires
the absolute value of its denominator.  Its sign does not affect
the conclusion that constant curvature gives a circle.}


\Needspace{5\baselineskip}
\par\bigskip
\noindent\textbf{22.4.6}\quad Show that requiring $J$, given by
\[
J=\int_{a}^{b}\!\!\int_{a}^{b}
K(x,t)\varphi(x)\varphi(t)\,dx\,dt,
\]
to have a stationary value subject to the normalizing condition
\[
\int_{a}^{b}\varphi^{2}(x)\,dx=1
\]
leads to a Hilbert--Schmidt integral equation of the form given
in Eq.~\erratum{(21.64)}, \erratum{provided the corresponding integral-operator
eigenvalue is nonzero. State separately what happens in the nullspace.}
\par\emph{Note.}~The kernel $K(x,t)$ is symmetric.
\erratum{Take real $\varphi\in L^2[a,b]$ and a real
$K\in L^2([a,b]^2)$, so the quadratic functional is well defined.}
\solution
Assume real functions and a real symmetric kernel,
$K(x,t)=K(t,x)$, on $[a,b]$, with integrals for which the following
variations are valid.  Define
\[
 (T\phi)(x)=\int_a^bK(x,t)\phi(t)\,dt,\qquad
 \mathcal A[\phi]=J[\phi]-\mu\left(\int_a^b\phi^2\,dx-1\right).
\]
Under $\phi\mapsto\phi+\epsilon\eta$,
\begin{align*}
 \left.\frac{dJ}{d\epsilon}\right|_0
 &=\int_a^b\int_a^bK(x,t)
       [\eta(x)\phi(t)+\phi(x)\eta(t)]\,dx\,dt\\
 &=2\int_a^b\eta(x)(T\phi)(x)\,dx.
\end{align*}
In the second term, interchange $x,t$ and use the symmetry of $K$.
The normalization contributes $2\int\phi\eta\,dx$.  Hence
\[
 \delta\mathcal A
 =2\int_a^b\eta(x)\bigl[(T\phi)(x)-\mu\phi(x)\bigr]\,dx=0.
\]
Arbitrariness of $\eta$ gives
\[
 \boxed{\int_a^bK(x,t)\phi(t)\,dt=\mu\phi(x)}.
\]
Multiplying by $\phi(x)$ and integrating shows that
$J=\mu\int\phi^2=\mu$.  For $\mu\ne0$, set $\lambda=1/\mu$:
\[
 \boxed{\phi(x)=\lambda\int_a^bK(x,t)\phi(t)\,dt}.
\]
\erratum{The textbook's reference to Eq.~(20.64) should be
Eq.~(21.64), the homogeneous Hilbert--Schmidt equation.
Moreover, the reciprocal form assumes $\mu\ne0$.  A normalized
element of the kernel of $T$ is also stationary, with $J=\mu=0$,
and cannot be represented by a finite reciprocal $\lambda$.}
For instance, $K=1$ on $[0,1]^2$ and
$\phi(x)=\sqrt{12}(x-\tfrac12)$ give
$\int\phi^2=1$, $\int\phi=0$, and $T\phi=0$.  The undivided eigenvalue
equation above includes this case.


\Needspace{5\baselineskip}
\par\bigskip
\noindent\textbf{22.4.7}\quad An unknown function satisfies the
differential equation
\[
y''+\left(\frac{\pi}{2}\right)^{2}y=0
\]
and the boundary conditions $\erratum{y'(0)=0}$, $y(1)=0$,
\erratum{with $y(0)=1$ fixing the normalization.}
\erratum{Regard this as the ground-state problem for
$-y''=\lambda y$, whose exact eigenvalue is $(\pi/2)^2$.}
\begin{enumerate}
\item[(a)] Calculate the approximation $\lambda=F[y_{\text{trial}}]$
for $y_{\text{trial}}=1-x^{2}$.
\item[(b)] Compare with the exact eigenvalue.
\end{enumerate}
\par\nobreak\hfill\emph{ANS.}\enspace
(a) $\lambda=2.5$.\quad(b) $\lambda/\lambda_{\text{exact}}=1.013$.
\solution
\erratum{The nonhomogeneous values $y(0)=1$, $y(1)=0$ alone do not
define a Sturm--Liouville eigenvalue problem with variable $\lambda$.
The intended Rayleigh problem has $y'(0)=0$, $y(1)=0$, with $y(0)=1$
fixing the scale of the eigenfunction.  The given exact solution
and both trial families in their optimum ranges obey that additional
Neumann condition.}
Indeed, the quotient to be used is
\[
 F[y]=\frac{\int_0^1y'^2\,dx}{\int_0^1y^2\,dx},\qquad y(1)=0.
\]
Varying its numerator minus $\lambda$ times its denominator gives
\[
 2[y'\eta]_0^1-2\int_0^1(y''+\lambda y)\eta\,dx=0.
\]
The endpoint $1$ has $\eta(1)=0$, whereas the free endpoint $0$
requires $y'(0)=0$.  The equation is $-y''=\lambda y$.

(a) For $y_{\rm trial}=1-x^2$,
\[
 \int_0^1y'^2\,dx=4\int_0^1x^2\,dx=\frac43,
\]
and
\[
 \int_0^1y^2\,dx
 =\int_0^1(1-2x^2+x^4)\,dx
 =1-\frac23+\frac15=\frac8{15}.
\]
Consequently
\[
 \boxed{\lambda_{\rm trial}=\frac{4/3}{8/15}=\frac52=2.5}.
\]

(b) The conditions $y'(0)=0$, $y(1)=0$ give
$y=C\cos(\sqrt\lambda x)$ and
$\cos\sqrt\lambda=0$.  The lowest eigenvalue and its normalized-at-zero
eigenfunction are
\[
 \lambda_{\rm exact}=\frac{\pi^2}{4},\qquad
 y_{\rm exact}=\cos\frac{\pi x}{2}.
\]
Thus
\[
 \boxed{\frac{\lambda_{\rm trial}}{\lambda_{\rm exact}}
       =\frac{10}{\pi^2}=1.013211836\ldots},
\]
an excess of approximately $1.32118\%$.


\Needspace{5\baselineskip}
\par\bigskip
\noindent\textbf{22.4.8}\quad In Exercise~22.4.7 use a trial
function $y=1-x^{n}$, \erratum{with $n>1/2$ for finite energy}.
\begin{enumerate}
\item[(a)] Find the value of $n$ that will minimize
$F[y_{\text{trial}}]$.
\item[(b)] Show that the optimum value of $n$ drives the ratio
$\lambda/\lambda_{\text{exact}}$ down to $1.003$.
\end{enumerate}
\par\nobreak\hfill\emph{ANS.}\enspace(a) $n=1.7247$.
\solution
Use the Rayleigh quotient and boundary interpretation established in
Exercise~22.4.7.  The trial $y=1-x^n$ has finite derivative energy
precisely for $n>\tfrac12$, because
\[
 N=\int_0^1y'^2\,dx
   =n^2\int_0^1x^{2n-2}\,dx=\frac{n^2}{2n-1}.
\]
Its squared norm is
\begin{align*}
 D&=\int_0^1(1-2x^n+x^{2n})\,dx\\
  &=1-\frac2{n+1}+\frac1{2n+1}\\
  &=\frac{(n+1)(2n+1)-2(2n+1)+(n+1)}
          {(n+1)(2n+1)}
   =\frac{2n^2}{(n+1)(2n+1)}.
\end{align*}
Therefore
\[
 F(n)=\frac{(n+1)(2n+1)}{2(2n-1)}.
\]
Differentiating without dropping the denominator derivative,
\[
 F'(n)
 =\frac{(4n+3)(4n-2)-4(2n^2+3n+1)}{(4n-2)^2}
 =\frac{4n^2-4n-5}{2(2n-1)^2}.
\]
The critical points solve $4n^2-4n-5=0$, so the admissible one is
\[
 \boxed{n_{\rm opt}=\frac{1+\sqrt6}{2}
       =1.724744871\ldots}.
\]
The other root is negative.  The derivative changes from negative to
positive at $n_{\rm opt}$, and $F(n)\to\infty$ at both ends of
$(\tfrac12,\infty)$, so this is the global minimum in the trial family.
Substitution yields
\[
 F_{\min}=\frac{5+2\sqrt6}{4},\qquad
 \boxed{\frac{F_{\min}}{\pi^2/4}
 =\frac{5+2\sqrt6}{\pi^2}=1.002976318\ldots}.
\]
The energy error is therefore approximately $0.297632\%$.

\erratum{The variational domain here is $n>\tfrac12$, not merely
$n>0$: smaller positive $n$ satisfy the endpoint values but have
infinite derivative energy.  The quadratic-form quotient
$\int y'^2/\int y^2$ remains the appropriate expression throughout
this domain; replacing it by $-\int yy''/\int y^2$ without checking
the boundary term is invalid for some of these trials.}
At the optimum $n>1$, $y'(0)=0$, so that boundary term vanishes.


\Needspace{5\baselineskip}
\par\bigskip
\noindent\textbf{22.4.9}\quad A quantum mechanical particle in a
sphere (Example~14.7.1) satisfies
\[
\nabla^{2}\psi+k^{2}\psi=0,
\]
with $k^{2}=2mE/\hbar^{2}$. The boundary condition is that
$\psi=0$ at $r=a$, where $a$ is the radius of the sphere. For
the ground state [where $\psi=\psi(r)$], try an approximate wave
function,
\[
\psi_{a}(r)=1-\left(\frac{r}{a}\right)^{2},
\]
and calculate an approximate eigenvalue $k_{a}^{2}$.
\erratum{Take $a>0$ and impose regularity at the origin.}
\par\emph{Hint.}~To determine $p(r)$ and $w(r)$, put your
equation in self-adjoint form (in spherical polar coordinates).
\par\nobreak\hfill\emph{ANS.}\enspace
$k_{a}^{2}=\dfrac{10.5}{a^{2}}$,\quad
$k_{\text{exact}}^{2}=\dfrac{\pi^{2}}{a^{2}}$.
\solution
For a radial function, the differential equation is
\[
 \frac1{r^2}\frac d{dr}(r^2\psi')+k^2\psi=0,
 \quad\hbox{or}\quad
 (r^2\psi')'+k^2r^2\psi=0.
\]
Thus the self-adjoint coefficients are
$p(r)=r^2$, $w(r)=r^2$, $q(r)=0$.  Multiplication by $\psi$ and
integration by parts gives
\[
 k^2=\frac{\int_0^ar^2\psi'^2\,dr}
           {\int_0^ar^2\psi^2\,dr}.
\]
The boundary term $[r^2\psi\psi']_0^a$ vanishes at $a$ because
$\psi(a)=0$, and at the origin for the regular radial solution.

For $\psi_{\rm trial}=1-r^2/a^2$,
\[
 N=\frac4{a^4}\int_0^ar^4\,dr=\frac{4a}{5},
\]
whereas
\begin{align*}
 D&=\int_0^a\left(r^2-\frac{2r^4}{a^2}+\frac{r^6}{a^4}\right)dr\\
  &=a^3\left(\frac13-\frac25+\frac17\right)
   =\frac{8a^3}{105}.
\end{align*}
Therefore
\[
 \boxed{k_{\rm trial}^2
 =\frac{4a/5}{8a^3/105}=\frac{21}{2a^2}
 =\frac{10.5}{a^2}}.
\]
The exact regular solution is proportional to
$\sin(kr)/(kr)$.  The smallest positive zero at $r=a$ is $ka=\pi$,
so
\[
 k_{\rm exact}^2=\frac{\pi^2}{a^2},\qquad
 \frac{k_{\rm trial}^2}{k_{\rm exact}^2}
 =\frac{21}{2\pi^2}=1.063872428\ldots.
\]
The resulting energy is approximately $6.38724\%$ above the exact
ground-state value, since $E=\hbar^2k^2/(2m)$.


\Needspace{5\baselineskip}
\par\bigskip
\noindent\textbf{22.4.10}\quad The wave equation for a quantum
mechanical oscillator may be written as
\[
\frac{d^{2}\psi(x)}{dx^{2}}+(\lambda-x^{2})\psi(x)=0,
\]
with $\lambda=1$ for the ground state; see Eq.~(18.17). Take
\[
\psi_{\text{trial}}=\begin{cases}
1-\dfrac{x^{2}}{a^{2}}, & x^{2}\le a^{2}\\[8pt]
0, & x^{2}>a^{2}
\end{cases}
\]
for the ground-state wave function (with $a^{2}$ an adjustable
parameter) and calculate the corresponding ground-state energy.
How much error do you have?
\erratum{Take $a>0$ and use the Rayleigh quotient with kinetic term
$\int|\psi'|^2dx$; the trial has derivative jumps at its cutoffs.}
\par\emph{Note.}~Your parabola is really not a very good
approximation to a Gaussian exponential. What improvements can
you suggest?
\solution
The dimensionless Hamiltonian is $H=-d^2/dx^2+x^2$, so its Rayleigh
quotient is
\[
 \lambda[\psi]
 =\frac{\int_{\mathbb R}(\psi'^2+x^2\psi^2)\,dx}
        {\int_{\mathbb R}\psi^2\,dx}.
\]
\erratum{The stated trial is continuous, but its derivative jumps at
$\pm a$.  It belongs to the quadratic-form domain, and the quotient
above is valid.  Treating its second derivative as an ordinary
function across those joins would omit delta-function terms.}
All three integrals reduce to $[-a,a]$:
\begin{align*}
 D&=2\int_0^a\left(1-\frac{2x^2}{a^2}+\frac{x^4}{a^4}\right)dx
   =2a\left(1-\frac23+\frac15\right)=\frac{16a}{15},\\
 N_{\rm kin}&=2\int_0^a\frac{4x^2}{a^4}\,dx=\frac8{3a},\\
 N_{\rm pot}&=2\int_0^a
  \left(x^2-\frac{2x^4}{a^2}+\frac{x^6}{a^4}\right)dx
  =2a^3\left(\frac13-\frac25+\frac17\right)=\frac{16a^3}{105}.
\end{align*}
Hence
\[
 \lambda(a)=\frac5{2a^2}+\frac{a^2}{7}.
\]
Put $A=a^2>0$.  Then
\[
 \frac{d\lambda}{dA}=-\frac5{2A^2}+\frac17=0
 \quad\Longrightarrow\quad
 A^2=\frac{35}{2},\qquad
 \boxed{a=\left(\frac{35}{2}\right)^{1/4}}.
\]
The second derivative $5/A^3$ is positive.  At the minimum the two
energy terms are equal:
\[
 \boxed{\lambda_{\min}=2\sqrt{\frac5{14}}
       =\sqrt{\frac{10}{7}}=1.195228609\ldots}.
\]
Thus the relative error is
\[
 \boxed{100(\lambda_{\min}-1)\%=19.5228609\ldots\%}.
\]
In conventional oscillator units,
$E_{\rm trial}=(\hbar\omega/2)\sqrt{10/7}$.
\erratum{The solution manual's final value $2\sqrt{10/7}$ has an
extra factor of two.  Substitution into its own expression
$5/(2a^2)+a^2/7$ gives $\sqrt{10/7}$.}

One improvement is to make both the value and slope vanish at the
cutoff:
\[
 \psi=(1-x^2/a^2)^2\quad (|x|\le a),\qquad \psi=0\quad(|x|>a).
\]
Writing $t=x/a$, the needed polynomial integrals are
\[
 \int_{-1}^1(1-t^2)^4dt=\frac{256}{315},\quad
 \int_{-1}^116t^2(1-t^2)^2dt=\frac{256}{105},\quad
 \int_{-1}^1t^2(1-t^2)^4dt=\frac{256}{3465}.
\]
They give
\[
 \lambda(a)=\frac3{a^2}+\frac{a^2}{11},\qquad
 a^4=33,\qquad
 \lambda_{\min}=2\sqrt{\frac3{11}}=1.044465936\ldots,
\]
only $4.44659\%$ high.  A still better family is the Gaussian
$\psi=e^{-\beta x^2/2}$, $\beta>0$.  Since
$\langle x^2\rangle=1/(2\beta)$ and
$\psi'=-\beta x\psi$,
\[
 \lambda(\beta)=\frac{\beta}{2}+\frac1{2\beta}\ge1,
\]
with equality at $\beta=1$, which recovers the exact ground state.


\Needspace{5\baselineskip}
\par\bigskip
\noindent\textbf{22.4.11}\quad The Schr\"odinger equation for a
central potential may be written as
\[
\mathcal{L}u(r)+\frac{\hbar^{2}l(l+1)}{2Mr^{2}}\,u(r)=Eu(r).
\]
The $l(l+1)$ term, the angular momentum barrier, comes from
splitting off the angular dependence. Compare Eq.~(9.80) (divide
that equation through by $-r^{2}$). Use the Rayleigh--Ritz
technique to show that $E>E_{0}$ \erratum{for $l>0$ and normalizable
bound states}, where $E_{0}$ is the \erratum{lowest} energy
eigenvalue of $\mathcal{L}u_{0}=E_{0}u_{0}$ corresponding to
$l=0$. This means that the ground state will have $l=0$, zero
angular momentum.
\erratum{Assume a self-adjoint radial Hamiltonian bounded below,
with the usual common quadratic-form domain.}
\par\emph{Hint.}~For a discrete orthonormal eigenbasis, expand
$u(r)=\erratum{\sum_{i=0}^{\infty}c_i u_i(r)}$, where
$\mathcal{L}u_i=E_i u_i$ and $E_i\ge E_0$.
\erratum{Do not require $c_0=1$; a spectral integral handles any
continuous part of the spectrum.}
\solution
\erratum{The strict inequality is for $l>0$ and a normalizable bound
eigenstate, with $E_0$ the lowest eigenvalue of the $l=0$ radial
operator.  For $l=0$ the ground state has equality.  A continuum
threshold need not have a strict gap above $E_0$.}
Normalize $u$ by $\int_0^\infty|u|^2dr=1$ and assume it belongs to the
quadratic-form domain of the radial Hamiltonian, including the
centrifugal term. For form-domain functions,
$\langle u,Lu\rangle$ below denotes the associated quadratic form;
it agrees with the operator inner product when $u$ is in the operator
domain. The variational lower bound for $L$ is
\[
 \langle u,Lu\rangle\ge E_0\langle u,u\rangle=E_0.
\]
For the discrete expansion in the hint, this follows directly from
$u=\sum_{i\ge0}c_i u_i$ in an orthonormal eigenbasis:
\[
 \sum_i|c_i|^2=1,\qquad
 \langle u,Lu\rangle
 =\sum_i|c_i|^2E_i
 =E_0+\sum_{i\ge1}|c_i|^2(E_i-E_0)\ge E_0.
\]
There is no need to assume $c_0\ne0$ or to divide by it.  If there
is continuous spectrum, the same argument uses its spectral
integral; the defining lower bound of $L$ suffices.

Multiply the radial equation by $u^*$ and integrate:
\[
 E=\langle u,Lu\rangle
   +\frac{\hbar^2l(l+1)}{2M}
      \int_0^\infty\frac{|u(r)|^2}{r^2}\,dr.
\]
For a nonzero normalizable state the last integral is strictly
positive: its integrand is nonnegative, and $r^{-2}>0$ at every
finite $r>0$.  For $l>0$ its coefficient is also positive.  Thus
\[
 \boxed{E>E_0\qquad(l>0)}.
\]
Every bound state with nonzero angular momentum lies above the
$l=0$ ground state.  Consequently, if a bound ground state exists
for the central-potential problem, it has $l=0$.
